% Options for packages loaded elsewhere
\PassOptionsToPackage{unicode}{hyperref}
\PassOptionsToPackage{hyphens}{url}
%
\documentclass[
  10pt,
]{article}
\usepackage{lmodern}
\usepackage{amssymb,amsmath}
\usepackage{ifxetex,ifluatex}
\ifnum 0\ifxetex 1\fi\ifluatex 1\fi=0 % if pdftex
  \usepackage[T1]{fontenc}
  \usepackage[utf8]{inputenc}
  \usepackage{textcomp} % provide euro and other symbols
\else % if luatex or xetex
  \usepackage{unicode-math}
  \defaultfontfeatures{Scale=MatchLowercase}
  \defaultfontfeatures[\rmfamily]{Ligatures=TeX,Scale=1}
\fi
% Use upquote if available, for straight quotes in verbatim environments
\IfFileExists{upquote.sty}{\usepackage{upquote}}{}
\IfFileExists{microtype.sty}{% use microtype if available
  \usepackage[]{microtype}
  \UseMicrotypeSet[protrusion]{basicmath} % disable protrusion for tt fonts
}{}
\makeatletter
\@ifundefined{KOMAClassName}{% if non-KOMA class
  \IfFileExists{parskip.sty}{%
    \usepackage{parskip}
  }{% else
    \setlength{\parindent}{0pt}
    \setlength{\parskip}{6pt plus 2pt minus 1pt}}
}{% if KOMA class
  \KOMAoptions{parskip=half}}
\makeatother
\usepackage{xcolor}
\IfFileExists{xurl.sty}{\usepackage{xurl}}{} % add URL line breaks if available
\IfFileExists{bookmark.sty}{\usepackage{bookmark}}{\usepackage{hyperref}}
\hypersetup{
  hidelinks,
  pdfcreator={LaTeX via pandoc}}
\urlstyle{same} % disable monospaced font for URLs
\usepackage[a4paper,margin=18mm]{geometry}
\setlength{\emergencystretch}{3em} % prevent overfull lines
\providecommand{\tightlist}{%
  \setlength{\itemsep}{0pt}\setlength{\parskip}{0pt}}
\setcounter{secnumdepth}{-\maxdimen} % remove section numbering

\usepackage{esint}
\usepackage{xeCJK}
\setmainfont{DejaVu Serif}
\setmonofont{DejaVu Sans Mono}
\setCJKmainfont{Noto Serif CJK SC}
\setCJKmonofont{Noto Sans Mono CJK SC}
\setmathfont{Asana Math}
\setmathfont{Latin Modern Math}[range=\lrcorner]
\setlength{\emergencystretch}{3em}
\allowdisplaybreaks
\hypersetup{hypertexnames=false}

\author{}
\date{}

\begin{document}

\hypertarget{linear-separation-and-actual-applications}{%
\section{Linear separation and actual
applications}\label{linear-separation-and-actual-applications}}

Research snapshot: 844e62f7304d9f890914a6a4370ec39291ddc26e.

Complete recent proofs: I40 closes linear order, I41 then removes scalar
variance and produces the actual BF extension. The historical open
statement at the end of the first paper's Section 4.3 is superseded by
the second paper's Theorem 1.

Read 00\_READ\_FIRST.md for superseded statements. Each source below is
complete and byte-preserved. Original local links remain relative to
their source; locate other texts using Volume 6. Formula numbers restart
in each source.

\begin{center}\rule{0.5\linewidth}{0.5pt}\end{center}

\hypertarget{source-1-20260927_mixed_current_and_signed_cores.md}{%
\subsection{Source 1:
20260927\_mixed\_current\_and\_signed\_cores.md}\label{source-1-20260927_mixed_current_and_signed_cores.md}}

Original path:
continuation/20260927\_mixed\_current\_and\_signed\_cores.md

SHA-256:
ea0c4955b56edf70fc5a19d5fd5f4553a1c0c8962a442cdd8c43362b6aabcf79

BEGIN COMPLETE SOURCE

\hypertarget{linear-einstein-separation-at-the-small-surgery-boundary}{%
\section{Linear Einstein separation at the small-surgery
boundary}\label{linear-einstein-separation-at-the-small-surgery-boundary}}

27 September 2026. Continues
\href{20260926_actual_topological_cost.md}{the actual-cost proof},
commit \texttt{168bea5}. The
\href{../guidance/source/20260927_web_critical_mixed_term.zh.md}{new web
contribution} is preserved verbatim. \textbf{The main result is proved
in Sections 4--5:} a refined same-spinor Hessian estimate closes the
beta=2 weak-Bianchi bootstrap, gives linear topological cost, and
determines the order of the relative Einstein defect on the previously
constructed small-h family. Sections 1--3 below incorporate the web
note's mixed-current estimate, model cancellation and curvature
identity, with explicit attribution; they are not claimed as newly
discovered here. Subsequent sections record further proof work. All
actual-solvability hypotheses remain those of the inherited functional
pencil. No conjectured linear gap is a premise. The zero-Weyl endpoint
and general Einstein realization remain separate problems.

For precision, the actual input on a closed connected oriented smooth
four-manifold (M,g) is a Spin\^{}c structure whose unperturbed
Seiberg--Witten equations are solvable for every metric conformal to g.
Its real Chern class a is fixed throughout the auxiliary weights; it may
be selected anew when the original metric changes. Put

\[
 p=(a_g^+)^2>0,\quad r=V^{-1}\int|\operatorname{Ric}+g|^2,\quad
 D_2=\frac32\int(s^2+12|W^+|^2)-72\pi^2p\ge0.
\]

Norms and the rescaled equations are those of the
\href{../mathematics/current/line1_functional_pencil.md}{functional
pencil}, Sections 1--3. There is no diameter, noncollapse, Sobolev,
injectivity-radius or full-curvature assumption. The new beta=2 proof
uses the supplied actual-field estimates, not an arbitrary approximately
closed form.

Write rho=\textbar Phi\textbar\^{}2, zeta=\textless Phi,nabla\_A
Phi\textgreater, with the Hermitian pairing linear in its second entry,
P=nabla\_A Phi, and Theta=(kappa+epsilon)V. Curvature two-form norms use
sum\_\{i\textless j\}\textbar F\_ij\textbar\_HS\^{}2. The selectable
smoothing error epsilon belongs to one actual weighted solution, not a
comparison of two solutions.

\hypertarget{the-full-rank-two-transgression-and-the-mixed-bound}{%
\subsection{1. The full rank-two transgression and the mixed
bound}\label{the-full-rank-two-transgression-and-the-mixed-bound}}

On \{rho\textgreater0\}, use u=Phi/sqrt(rho) and a local unitary
complement v. In this frame let the connection matrix and the connection
fixing u be

\[
 \mathcal A=\begin{pmatrix}a&-\bar b\\b&d\end{pmatrix},
 \qquad \mathcal A^0=\begin{pmatrix}0&0\\0&d\end{pmatrix}.
\]

Here a is an imaginary one-form, not the topological class. Set
B\_c=mathcal A-mathcal A\^{}0 and F\^{}t the curvature of mathcal
A\^{}0+tB\_c. The global transgression on the nonzero set is

\[
 T_u=\frac1{4\pi^2}\int_0^1
 [\operatorname{tr}(B_c\wedge F^t)
       -\operatorname{tr}B_c\wedge\operatorname{tr}F^t]dt,
 \qquad dT_u=c_2(F).
\]

The web note's exact expansion is

\[
 4\pi^2T_u=-a\wedge F_{22}
 +\tfrac12(b\wedge F_{12}-\bar b\wedge F_{21})
 +a\wedge\bar b\wedge b.                         \tag{1.1}
\]

Indeed F\^{}0\_22=F\_22-bar b wedge b, and
F\^{}t=(1-t)F\^{}0+tF-t(1-t)B\_c wedge B\_c; substitution and
integration in t give (1.1). This calculation retains the matrix order
and wedge signs. The expression is invariant under changing v by a unit
complex scalar.

For beta=2 put T=\{1/2\textless rho\textless1\}. On T,

\[
 \zeta=\tfrac12d\rho+\rho a,\quad
 |d\rho|\le2|\zeta|,\quad |a|\le2|\zeta|,
 \quad |b|\le\sqrt2|P|.                             \tag{1.2}
\]

Choose chi=(1-eta)\^{}2 with eta=0 below 1/2, eta=1 above 1 and
\textbar eta'\textbar\textless=4. The inherited signed identity and the
zero-set extension give

\[
 c_2(S^+)[M]=\int\chi c_2(F)+\int d\chi\wedge T_u,
 \quad
 16\pi^2c_2(F)=
 (|F_A^+|^2-|F_A^-|^2-2|\Omega^+|^2+2|\Omega^-|^2)d\mu.
                                                               \tag{1.3}
\]

Let G\^{}2=\textbar Omega\textbar\^{}2+\textbar F\_A\^{}+\textbar\^{}2.
The actual estimates already give
int\_\{rho\textless=1\}G\^{}2\textless=CTheta and the same bound for the
positive part in (1.3). Since off-diagonal curvature has no central
part, (1.1) yields

\[
 |d\chi\wedge T_u|\le C|\chi'|
 (G|\zeta||P|+|F_A^-||\zeta|^2+|\zeta|^2|P|^2).
\]

Young's inequality controls the central term by
chi\textbar F\_A\^{}-\textbar\^{}2/(32pi\^{}2)+C(chi')\^{}2\textbar zeta\textbar\^{}4/chi.
The first summand is absorbed by the negative term in (1.3). The ratio
(chi')\^{}2/chi\textless=64 and
\textbar zeta\textbar\^{}2\textless=rho\textbar P\textbar\^{}2\textless=\textbar P\textbar\^{}2
on T control the second. Consequently:

\textbf{Proposition 1 (web mixed-current refinement).}

\[
 \boxed{c_2(S^+)[M]\le C\Theta+C\int_T|\zeta|^2|P|^2.} \tag{1.4}
\]

In the exact small-h family of the preceding continuation, any E=o(h)
therefore requires
int\_T\textbar zeta\textbar\^{}2\textbar P\textbar\^{}2\textgreater=cAh.
This is stronger than the coarser quartic concentration criterion, but
does not exclude that concentration.

There is a further exact simplification useful for retaining the sign.
Since F\_11=da-bar b wedge b and tr F=F\_A,

\[
 \boxed{4\pi^2T_u=a\wedge da-a\wedge F_A
 +\tfrac12(b\wedge F_{12}-\bar b\wedge F_{21}).}       \tag{1.5}
\]

In particular the two diagonal terms in the cutoff flux are exactly

\[
 \frac{\chi'(\rho)}{\rho^2}d\rho\wedge\zeta\wedge d\zeta
 -\frac{\chi'(\rho)}\rho d\rho\wedge\zeta\wedge F_A.  \tag{1.6}
\]

The real exact part of zeta drops out after wedging with d rho. Formula
(1.6) is signed; its first term is not asserted nonpositive. Replacing
it by its absolute value loses the current's phase/radial correlation.

\hypertarget{axial-model-cancellation-in-the-actual-torsion-equation}{%
\subsection{2. Axial model cancellation in the actual torsion
equation}\label{axial-model-cancellation-in-the-actual-torsion-equation}}

Keep the globally smooth cubic field

\[
 t_\Phi=\sum_i c(e^i)\Phi\otimes(\Phi\wedge P_i)\in S^-\otimes L.
\]

The preceding exact calculation gives D\^{}-\emph{\{3A\}t\_Phi=-4 sum\_i
P\_i tensor(Phi wedge P\_i)+mathcal K(W\^{}+,Phi). For the actual axial
model W\_0=A}\{omega,1/beta\}, the web note observes

\[
 \mathcal K(W_0,\Phi)=0.                              \tag{2.1}
\]

To verify this, at Phi!=0 choose J\_1 Phi=iPhi. Its wedge factor is
zero. The other model eigenvalues coincide and J\_3 Phi=+/-iJ\_2 Phi, so
the other tensor factors cancel by 1+i\^{}2=0. At Phi=0 the polynomial
vanishes. Thus throughout the actual finite-parameter range

\[
 D^-_{3A}t_\Phi=-4\sum_i P_i\otimes(\Phi\wedge P_i)
       +\mathcal K(W^+-W_0,\Phi).                     \tag{2.2}
\]

With the inherited uniform Hessian budget (in particular
a\^{}2\textgreater=T\_+), and kappa\textless=1 with a sufficiently small
chosen epsilon,

\[
 \|D^-_{3A}t_\Phi\|_{4/3,\mathrm{avg}}
       \le C\sqrt{\kappa+\varepsilon},\qquad
 \|D^-_{3A}t_\Phi\|_{2,\mathrm{avg}}^2
       \le C\left(\kappa+\varepsilon+\fint|\zeta|^2|P|^2\right).
                                                               \tag{2.3}
\]

Unlike the previous coarse estimate, (2.3) has no beta\^{}\{-1\} term.
It does not assert nonzero torsion at the KE endpoint.

The mass calculation also has an all-parameter form. Write
q\_beta=4(1+1/beta), rho\_beta=2sqrt(q\_beta)/3 and
U\_beta=(4/3)(2/beta-1). The previous identities give, uniformly,

\[
 \boxed{\fint|t_\Phi|^2
 =\frac{16q_\beta^{3/2}}{81}(1-2/\beta)
       +O(\sqrt{\kappa+\varepsilon}).}                \tag{2.4}
\]

Indeed \textbar t\textbar\^{}2 differs from (rho/2)\textbar nabla
psi\textbar\^{}2 by an integrable error O(Theta); replacing rho by
rho\_beta costs O(sqrt(kappa+epsilon)) in the average, using the L4
gradient budget. Also avg\textbar nabla psi\textbar\^{}2=4rho\_beta
K\_A/V+O(sqrt(kappa+epsilon)), and K\_A/V=-U\_beta
rho\_beta/4+O(sqrt(kappa+epsilon)). Their product is -U\_beta
rho\_beta\^{}3/2, which is (2.4). The zero-end value is 128/81, whereas
the exact beta=2 leading term is zero.

\hypertarget{the-additional-actual-curvature-equation-and-its-sign}{%
\subsection{3. The additional actual curvature equation and its
sign}\label{the-additional-actual-curvature-equation-and-its-sign}}

The web note supplies a further equation, not an ordinary new monopole.
Use the imaginary-valued quadratic map q\_- such that
c(q\_-(t))t=(\textbar t\textbar\^{}2/2)t. Then

\[
 \rho^3F_A^-+2q_-(t_\Phi)=2\rho^2(d\zeta)^-+\mathcal R,
 \quad
 |\mathcal R|\le C\left(\rho^3|\operatorname{Ric}^\circ|
                  +\rho^{3/2}|\zeta||P|\right).       \tag{3.1}
\]

Here is a pointwise normalization check and derivation. On
rho\textgreater0 write Phi=sqrt(rho)u, P=alpha u+beta v,
alpha=zeta/sqrt(rho). Decompose beta=beta\_N+beta\_L using the kernel of
c(.)v. The Dirac equation gives
\textbar beta\_L\textbar\textless=\textbar alpha\textbar. In the
Clifford matrices of the previous proof this kernel consists of
beta\_N=(iz,z,iw,w), with t=2i rho(z,-w) on the kernel part. Directly,
q\_-(t)=rho\^{}2(bar beta\_N wedge beta\_N)\^{}-.

Now use (d
zeta)\_ij=\textless P\_i,P\_j\textgreater-\textless P\_j,P\_i\textgreater+\textless Phi,F\^{}S+\_ij
Phi\textgreater. The anti-self-dual central curvature contributes rho
F\_A\^{}-/2 and the traceless contribution is bounded by C
rho\textbar Ric\^{}circ\textbar. All terms in bar beta wedge beta-bar
beta\_N wedge beta\_N, and in bar alpha wedge alpha, are bounded by
C\textbar alpha\textbar\textbar P\textbar. Multiplication by 2rho\^{}2
gives (3.1). Both sides extend polynomially through zeros; the bound
extends by continuity.

For completeness the same scalar-potential spinor equation also gives

\[
 \delta\zeta=-|P|^2-U\rho/4,\qquad
 \delta\operatorname{Im}\zeta=0,\qquad
 \Delta\rho=-2|P|^2-U\rho/2.                         \tag{3.2}
\]

The Laplacian convention is Delta=d\emph{d\textgreater=0. The first
equation follows by differentiating \textless Phi,P\textgreater{} and
using nabla} nabla Phi=-U Phi/4. These are additional actual equations,
not independent free fields.

The derivative term in (3.1) can be written rho\^{}2 d zeta=d(rho\^{}2
zeta)-2rho d rho wedge zeta. Its first part is the derivative of a small
L2 field, the second is small in L1. This does not justify testing by an
uncontrolled derivative of q\_-(t).

Even a vanishing L2 Dirac residual would not finish the argument:

\[
 \int|D^-_{3A}t|^2=\int|\nabla t|^2+\int(s/4)|t|^2
       +\tfrac32\int\langle c(F_A^-)t,t\rangle.       \tag{3.3}
\]

The scalar term is negative here. The sign in (3.1) is opposite the
ordinary monopole sign. In particular no positive vanishing theorem is
being invoked. Actual U(1)-twisted Dirac modes on Eguchi--Hanson space
are also known; see
\href{https://arxiv.org/abs/2309.08453}{Franchetti--Krasnov}.

\hypertarget{removing-the-background-volume-from-the-second-order-budget}{%
\subsection{4. Removing the background volume from the second-order
budget}\label{removing-the-background-volume-from-the-second-order-budget}}

This is the decisive new step. The previous coarse Hessian estimate paid
O(V) for all background curvature. At beta=2 the bounded KE curvature
model annihilates the distinguished spinor at its calibrated amplitude.
Keeping this fact in the exact Hessian identity, and then using the
existing weak-Bianchi self-test again, closes the remaining estimate.

We distinguish the spinor gradient P\_i from the scalar Ricci energy
P\_Ric=int\textbar Ric\^{}circ\textbar\^{}2. Set

\[
 T_+=2\chi+3\tau,\quad p=(a_g^+)^2,\quad
 K=\int|\nabla_A\Phi|^2,\quad B=\int|F_A^-|^2,
 \quad \Xi=rV+D_2,\quad\Theta=\Xi+\varepsilon V.
\]

The actual beta=2 supplied solution satisfies, with universal constants,

\[
 |\Phi|\le2,\quad \rho_0=4/\sqrt6,\quad c=\sqrt6,
\] \[
 \int\bigl((\rho-\rho_0)^2+(f-c)^2+U^2+|d\rho|^2
       +|\operatorname{Im}\zeta|^2+|W^+-W_0|^2
       +|\operatorname{Ric}+g|^2\bigr)\le C\Theta,
 \quad J_+\le C\Theta.                              \tag{4.1}
\]

Here U=s+f rho, W\_0=A\_\{omega,1/2\}, omega=sqrt(3)psi/2. These are the
inherited estimates, with their exact same-solution quantifiers,
collected in Sections 1--3 of the preceding actual-cost proof. Its
localized equation and remainder identity also give

\[
 \int_{\{\rho\le1\}}|\nabla_A\Phi|^2\le C\Theta,
 \qquad |8\sqrt6 K-\Sigma|\le C\Theta.               \tag{4.2}
\]

\hypertarget{a-refined-hessian-lemma}{%
\subsubsection{4.1 A refined Hessian
lemma}\label{a-refined-hessian-lemma}}

\textbf{Lemma 2.} The same solution obeys

\[
 \boxed{\int\bigl(|\nabla_A^2\Phi|^2+|\nabla_A\Phi|^4\bigr)
               \le C(K+\Theta+B).}                  \tag{4.3}
\]

\textbf{Proof.} Form an algebraic positive-spinor curvature F\_b using
s\_b=-4, Ric\_b\^{}circ=0, W\_b\^{}+=W\_0, F\_\{A,b\}\^{}+=i c
sigma(Phi), and F\_\{A,b\}\^{}-=0. It need not itself be the curvature
of a connection; it is used only algebraically inside the actual Hessian
identity. Since Phi is bounded, \textbar F\_b\textbar\textless=C, and
(4.1) gives

\[
 \|F^{S^+}-F_b\|_2^2\le C(\Theta+B).                \tag{4.4}
\]

The more important bound is

\[
                  |F_b\Phi|^2\le C(\rho-\rho_0)^2.  \tag{4.5}
\]

For an explicit check choose the self-dual basis along psi. W\_0 has
eigenvalues -rho\^{}2/4,rho\^{}2/8,rho\^{}2/8. With the
Clifford/curvature conventions of the supplied equations, the axis and
transverse coefficients of F\_b Phi, up to their common fixed
normalizations, are

\[
 \rho^2/4+1/3-c\rho/4
     =\tfrac14(\rho-4/\sqrt6)(\rho-2/\sqrt6),
 \qquad
 \rho^2/8-1/3
     =\tfrac18(\rho-4/\sqrt6)(\rho+4/\sqrt6).
\]

They vanish at rho\_0; bounded rho proves (4.5), including zeros.
Equivalently one can split off the unit-amplitude KE model on the
nonzero set, which annihilates Phi exactly. No derivative of either
model is used in this step.

Put H=nabla\_A\^{}2 Phi and L=tr H. The exact covariant identity is

\[
 \|H\|_2^2=\|L\|_2^2
 -\int\operatorname{Ric}_{ij}\operatorname{Re}\langle P_i,P_j\rangle
 +\sum_{ij}\int\operatorname{Re}\langle F_{ij}^{S^+}P_j,P_i\rangle
 +\tfrac12\sum_{ij}\|F_{ij}^{S^+}\Phi\|_2^2.         \tag{4.6}
\]

The sums in (4.6) are double sums; their norms differ by a fixed factor
from the exterior-curvature norm used elsewhere. The identity follows by
interchanging Hessian slots before commuting on P. It differentiates
neither curvature nor f.~Bounded-section integration by parts also gives
\textbar\textbar P\textbar\textbar\_4\^{}2\textless=8\textbar\textbar H\textbar\textbar\_2.
Now \textbar\textbar L\textbar\textbar\_2\^{}2\textless=CTheta by the
actual spinor equation. Split Ric=-g+(Ric+g), and F=F\_b+(F-F\_b). The
bounded background terms cost CK, the error terms cost C
sqrt(Theta+B)\textbar\textbar H\textbar\textbar\_2, and (4.4)--(4.5)
give \textbar\textbar F Phi\textbar\textbar\_2\^{}2\textless=C(Theta+B).
Thus

\[
 \|H\|_2^2\le C(K+\Theta+B)
                 +C\sqrt{\Theta+B}\,\|H\|_2.
\]

Young's inequality and the bounded-section interpolation prove (4.3).
There is no additive V term. QED.

\hypertarget{the-square-condition-absorbs-the-remaining-determinant-energy}{%
\subsubsection{4.2 The square condition absorbs the remaining
determinant
energy}\label{the-square-condition-absorbs-the-remaining-determinant-energy}}

Assume a\^{}2\textgreater=T\_+. The determinant and Chern--Weil
identities, with all constants retained, are

\[
 B=4\pi^2(p-a^2)+J_+,
 \qquad
 36\pi^2(p-T_+)=\tfrac92P_{\rm Ric}+\tfrac98\Sigma-D_2.
                                                               \tag{4.7}
\]

Consequently

\[
 B\le\tfrac12P_{\rm Ric}+\tfrac18\Sigma-D_2/9+J_+
       \le\sqrt6 K+C\Theta.                          \tag{4.8}
\]

This is not a smallness assumption on B: it bounds B by the unknown
spinor energy that the next self-test will absorb. We used
P\_Ric\textless=rV, D\_2\textgreater=0 and (4.1)--(4.2). Thus (4.3)
improves to

\[
             \|H\|_2^2+\|P\|_4^4\le C(K+\Theta).   \tag{4.9}
\]

\hypertarget{closing-the-weak-bianchi-bootstrap}{%
\subsubsection{4.3 Closing the weak-Bianchi
bootstrap}\label{closing-the-weak-bianchi-bootstrap}}

Let J=div W(W\_0), where W(W\_0) is the Weyl four-tensor corresponding
to the smooth polynomial W\_0. The inherited weak contracted Bianchi
identity gives

\[
 \|J\|_2^2\le C\sqrt\Theta\,\|\nabla J\|_2.
\]

W\_0 is a parallel polynomial of degree four in bounded Phi.
Consequently \textbar nabla J\textbar\textless=C\textbar nabla\^{}2
W\_0\textbar{} and (4.9) give

\[
                    \|J\|_2^2\le C\sqrt{\Theta(K+\Theta)}.
                                                               \tag{4.10}
\]

No commutation of these derivatives, and hence no new full-curvature
term, is needed. The same actual J is the test field.

The pointwise polynomial divergence identity from the existing
\href{../mathematics/current/round21_analytic.md}{bounded-form
criterion}, Section 6, with t=1/2 and
v=\textbar omega\textbar\^{}2=3rho\^{}2/4, gives

\[
 \tfrac14 v(|\nabla\omega|^2+|d\omega|^2)
       \le |J|^2+\tfrac13|dv|^2.                     \tag{4.11}
\]

On \{rho\textgreater1\}, use \textbar nabla
omega\textbar\^{}2=3rho\textbar P\textbar\^{}2-3\textbar Im
zeta\textbar\^{}2 and dv=(3/2)rho d rho. Equations (4.1), (4.2), (4.11)
yield

\[
                       K\le C(\Theta+\|J\|_2^2).    \tag{4.12}
\]

The low-amplitude integral in (4.2) is essential: we did not divide by
rho at its zeros. Combining (4.10)--(4.12) and applying Young once more,

\[
 K\le C\Theta+C\sqrt{\Theta K}\quad\Longrightarrow\quad
                         \boxed{K\le C\Theta.}       \tag{4.13}
\]

This proves the mixed-term smallness sought in the web brief as a
consequence of the original self-test with a sharper budget, rather than
an assumed regularity theorem for the mixed term itself.

\textbf{Theorem 3 (linear beta=2 cost).} Under the inherited actual
conformal monopole hypotheses, p\textgreater0 and
a\^{}2\textgreater=T\_+, every normalized metric with
kappa=r+D\_2/V\textless=1 has supplied smooth solutions with any
sufficiently small chosen allowance epsilon, for which

\[
 K+B+\int(|\nabla_A^2\Phi|^2+|\nabla_A\Phi|^4)
             +\int|\zeta|^2|\nabla_A\Phi|^2\le C\Theta. 
                                                               \tag{4.14}
\]

Moreover, the solution-independent quantities satisfy

\[
 \boxed{0\le\Sigma\le C(rV+D_2),\qquad
         0\le p-T_+\le C_T(rV+D_2).}                 \tag{4.15}
\]

The nonnegative scalar slack follows from the classical scalar estimate;
p\textgreater=a\^{}2\textgreater=T\_+. Bounds (4.8), (4.9), (4.13) and
\textbar zeta\textbar\textless=2\textbar P\textbar{} prove (4.14).
Equations (4.2), (4.7) prove (4.15) after letting epsilon tend to zero
on the fixed metric and class. No convergence of the actual solutions is
asserted or needed. In particular,

\[
                     c_2(S^+)[M]\le C(rV+D_2).       \tag{4.16}
\]

For the transition set, (4.14) gives
int\_T\textbar zeta\textbar\^{}2\textbar P\textbar\^{}2\textless=CTheta.
Choosing epsilon=o(h) proves the web brief's required o(hV) bound when
kappa=o(h). This now contradicts the positive-charge necessity (1.4).

The theorem controls rV, including scalar variance and normalization
error; it does \textbf{not} assert the stronger unproved inequality
Sigma\textless=C(P\_Ric+D\_2).

A square-cost variant also follows from the same proof. Retain Theorem
3's hypotheses except a\^{}2\textgreater=T\_+. If Y=(T\_+-a\^{}2)\_+,
replace (4.8) by B\textless=sqrt6 K+CTheta+4pi\^{}2Y. Then

\[
 K+\Sigma+(p-T_+)_+\le C(\Theta+\sqrt{\Theta Y}),
 \quad
 B+\|\nabla_A^2\Phi\|_2^2+\|\nabla_A\Phi\|_4^4
                   \le C(\Theta+Y).                 \tag{4.17}
\]

The mixed term is controlled by the second, not the first, bound in
(4.17). Thus a square shortfall remains a genuine cost. No such cost
appears in (4.14)--(4.16).

\hypertarget{transfer-to-parameters-near-the-ke-endpoint}{%
\subsubsection{4.4 Transfer to parameters near the KE
endpoint}\label{transfer-to-parameters-near-the-ke-endpoint}}

This transfer uses no new field or gauge hypothesis. For
2\textless=beta\textless=3 let delta=beta-2 and Xi\_beta=rV+D\_beta. The
exact deficit decomposition D\_beta=A\_0+beta\^{}(-1)int z\^{}2, where
z=s+beta w and A\_0=Delta\_++int q\_-\^{}2\textgreater=0, gives

\[
 D_2=A_0+\tfrac12\int(z-\delta w)^2
       \le\beta D_\beta+\delta^2\int w^2.
\]

Since w=(z-(s+4)+4)/beta, two applications of the squared triangle
inequality give

\[
 \int w^2\le\frac4\beta D_\beta+
                \frac{16}{\beta^2}rV+\frac{32}{\beta^2}V
             \le2D_\beta+4rV+8V.
\]

Consequently

\[
                    \Xi_2\le5\Xi_\beta+8(\beta-2)^2V. \tag{4.18}
\]

Whenever kappa\_2\textless=1, Theorem 3 controls all its spinor
quantities by C{[}Xi\_beta+(beta-2)\^{}2V+epsilon V{]}. The majorant
f\_0=max(sqrt(max(-s+w,3)),(-s)\_+/4) does not depend on beta, so the
same smooth majorants and actual solutions serve this comparison. This
is a nearby-parameter corollary, not an estimate at
beta-\textgreater infinity.

\hypertarget{a-linear-two-chirality-gap-and-the-optimal-small-h-order}{%
\subsection{5. A linear two-chirality gap and the optimal small-h
order}\label{a-linear-two-chirality-gap-and-the-optimal-small-h-order}}

Assume that for every smooth metric on M one can choose a single actual
conformally solvable Spin\^{}c structure with

\[
 p\ge A>0,\quad a^2\ge T_+=A(1-h),\quad 0<h\le1/3.
\]

As before put

\[
 U=\frac{3\min(T_+,T_-)}A,\quad
 e=\frac{R}{8\pi^2A},\quad R=\int|\operatorname{Ric}|^2,
 \quad\mathcal E=\inf_{\lambda\ne0}
       \frac{\int|\operatorname{Ric}-\lambda g|^2}{\lambda^2V}.
\]

\textbf{Theorem 4 (linear topological window).} There is a universal
c\textgreater0 such that every smooth metric under these hypotheses
satisfies

\[
              \boxed{\mathcal E+(U-1-h)_+\ge c h.}    \tag{5.1}
\]

\textbf{Proof.} R\textgreater0 by the positive actual scalar budget. If
U\textless=0, e(1-E)\textless=U gives E\textgreater=1 and the conclusion
is immediate. Otherwise U\textless=3. Set Delta=(U-1-h)\_+. If
E+Delta\textgreater=1/49, then E+Delta\textgreater=(3/49)h. It remains
to consider the smaller branch, in particular E\textless1/2.

Let I=int s. If I\textgreater=0, the negative scalar estimate gives
32pi\^{}2A\textless=int s\_-\^{}2\textless=L\_s\textless=RE, so
4\textless=eE. Combining with e(1-E)\textless=U gives
E\textgreater=4/(U+4)\textgreater=4/7, a contradiction. Thus
I\textless0. Rescale by -R/I so that r=E and V=R/(4-E); p,D\_2,R,P\_Ric
are unchanged. Write z=p/A-1\textgreater=0. The exact identity is

\[
          D_2/A=36\pi^2[e-(1+h)-2z].                 \tag{5.2}
\]

Hence e\textgreater=1+h, e\textless=U/(1-E)\textless=6, and

\[
 D_2/A\le72\pi^2\Delta+96\pi^2E,\quad
 rV/A=\frac{8\pi^2eE}{4-E}\le\frac{96}{7}\pi^2E.
\]

Also A/V=(4-E)/(8pi\^{}2e)\textless=1/(2pi\^{}2). Thus

\[
 \Xi/A\le112\pi^2(E+\Delta),\qquad
 \kappa\le49(E+\Delta)<1.                            \tag{5.3}
\]

The second bound in (5.3) follows directly from r=E and
D\_2/V\textless=36Delta+48E, not from the coarser first bound alone. Now
p-T\_+\textgreater=Ah and Theorem 3 give h\textless=C\_T
Xi/A\textless=112pi\^{}2 C\_T(E+Delta). Take c=min(3/49,1/(112pi\^{}2
C\_T)). This is independent of topology, volume, class choice and
metric. QED.

This uses the same chiral upper budget as the preceding h\^{}2 theorem.
Its new content comes from (4.15), not from a new Chern--Weil identity.
The old optimized F(h) inequality remains valid; (5.1) supplies a
linear-width exclusion near the classical Ricci-energy boundary 1+h. It
is not a claim that (5.1) dominates every previous inequality at every
parameter value.

\textbf{Corollary 5 (sharp order on the realized family).} On the exact
complex-hyperbolic-cover blowup/handle family constructed in Section 3,
Theorem 5 of \href{20260926_boundary_compatibility.md}{the boundary
continuation},

\[
 \boxed{c_1h_j\le\inf_g\mathcal E(g)
                   \le\frac{4h_j}{1+h_j}\le4h_j}
                  \qquad\text{for all sufficiently large }j. \tag{5.4}
\]

Indeed U\_j-F(h\_j)=o(h\_j\^{}2), and
F(h)-1-h=4h\^{}2/(1+sqrt(1-3h))\^{}2. Therefore
(U\_j-1-h\_j)\_+=O(h\_j\^{}2); absorb it into ch\_j/2 in (5.1). The
upper bound is the previously established consequence of LeBrun's actual
shrinking-core construction, not a new gluing theorem. Thus the optimal
\textbf{order} on this family is linear. The optimal leading constant
over all metrics has not been determined.

Combining E\textless=4delta\_sigma\^{}2 with (5.4) yields
delta\_sigma\textgreater=c sqrt(h) on this family for either sign. This
is a pointwise-input obstruction, not a pointwise near-Einstein
construction or a full solution of Yau26.

There is also a more informative feasible-frontier consequence. On the
optimally negative-normalized branch with kappa\textless=1, set
alpha=(36pi\^{}2 C\_T)\^{}(-1). From p-T\_+=A(h+z) and (4.15),

\[
 e\left(1+\frac{2\mathcal E}{9(4-\mathcal E)}\right)
 \ge 1+(1+\alpha)h+(2+\alpha)z,\qquad z=p/A-1\ge0.
                                                               \tag{5.5}
\]

For E\textless1, combining with e(1-E)\textless=U gives

\[
 U\ge\frac{(1-\mathcal E)[1+(1+\alpha)h]}
                 {1+2\mathcal E/[9(4-\mathcal E)]}.  \tag{5.6}
\]

These formulas use precisely the same hypotheses as the small-kappa
branch, not arbitrary metrics with uncontrolled D\_2/V. In particular,
at an exact negative Einstein metric in that branch they force
U\textgreater=1+(1+alpha)h. The all-metric statement remains (5.1),
whose separate large-defect branch supplies the needed global
quantifier.

\hypertarget{relaxing-the-class-square-hypothesis-without-losing-the-order}{%
\subsubsection{5.1 Relaxing the class-square hypothesis without losing
the
order}\label{relaxing-the-class-square-hypothesis-without-losing-the-order}}

The square-cost estimate (4.17) gives a genuine scope extension. Retain
the hypotheses of Theorem 4 except a\^{}2\textgreater=T\_+, and, for the
actual class selected at g, put y=(T\_+-a\^{}2)\_+/A. Then

\[
 \boxed{\mathcal E(g)+(U-1-h)_+\ge c\frac{h^2}{h+y}.} \tag{5.7}
\]

To prove it, the negative-average and rescaling arguments of Theorem 4
use p\textgreater=A, not the class-square condition. In the small branch
let x=Xi/A. Since p-T\_+\textgreater=Ah, (4.17), with epsilon sent to
zero, gives h\textless=C(x+sqrt(xy)). If x\textless=h, squaring gives
h\^{}2\textless=2C\^{}2 x(x+y)\textless=2C\^{}2 x(h+y). If
x\textgreater h, then x\textgreater h\textgreater=h\^{}2/(h+y). Combine
with x\textless=112pi\^{}2(E+Delta). In the large branch use
h\^{}2/(h+y)\textless=h\textless=1/3. This proves (5.7) with a universal
positive constant.

In particular a metricwise actual supply with
(T\_+-a\^{}2)\_+\textless=L A h, for a fixed L\textgreater=0, still
gives

\[
               \mathcal E(g)+(U-1-h)_+\ge c h/(1+L). \tag{5.8}
\]

This replaces the exact square inequality by a controlled shortfall at
the surgery scale. It does not assert the existence of any additional
gauge supplier. The realized family of Corollary 5 has y=0.

\hypertarget{construction-side-scope-inverse-surgery-curvature-costs}{%
\subsection{6. Construction-side scope: inverse-surgery curvature
costs}\label{construction-side-scope-inverse-surgery-curvature-costs}}

The parallel construction analysis gives a separate geometric
explanation of why local insertion alone cannot beat the linear order.
It is not used to prove Theorem 3.

Let N be a common oriented exterior in M and X, where T\_+(X)=A and
T\_+(M)=A(1-h). Suppose a smooth metric ghat on X agrees with g on N.
Write C\_X=X\textbackslash N, C\_M=M\textbackslash N and
H\_g=s\_g\^{}2/24+2\textbar W\_g\^{}+\textbar\^{}2. Global Chern--Weil
on the two closed manifolds gives the exact identity

\[
 P_M=P_X+2\int_{C_M}H_g-2\int_{C_X}H_{\widehat g}
                              +8\pi^2Ah.             \tag{6.1}
\]

Every term other than the cap energy has the displayed favorable sign.
Let gamma=(4pi\^{}2A)\^{}\{-1\}int\_\{C\_X\}H\_ghat and b=4(h-gamma)\_+.
For U\textgreater0 the chiral upper bound consequently gives

\[
                         \mathcal E(g)\ge\frac b{U+b}.
                                                               \tag{6.2}
\]

For E\textless1 this follows from E\textgreater=4P\_M/R\textgreater=b/e
and e(1-E)\textless=U; E\textgreater=1 is immediate. No closeness to a
complex-hyperbolic exterior is assumed. Thus low-energy inverse-surgery
caps with energy o(Ah) would force liminf E/h\textgreater=4 when
h-\textgreater0 and U-\textgreater1. Conversely every such replacement
of a hypothetical sublinear metric would have to cost at least
4pi\^{}2Ah-o(Ah). The main theorem now excludes that sublinear case in
the stated gauge/window family independently of the existence of caps.

Here is an explicit sufficient geometric criterion for cheap caps. On
disjoint spherical collars suppose, in coordinates at arbitrary scales
ell\_i, g=ell\_i\^{}2(delta+q\_i), with
\textbar\textbar q\_i\textbar\textbar\_\{C\^{}2(1\textless\textbar x\textbar\textless2)\}
\textless=epsilon\_i\textless=epsilon\_0. A fixed radial cutoff which is
zero for r\textless=5/4 and one for r\textgreater=7/4 extends the metric
to a Euclidean inner ball, unchanged at the exterior. The scaled
curvature is O(epsilon\_i), volume is bounded, and dimension-four
scaling cancels; hence

\[
                         \int_{\mathrm{cap}_i}H\le C\epsilon_i^2.
                                                               \tag{6.3}
\]

Thus sum epsilon\_i\^{}2=o(Ah) suffices, even with unbounded numbers of
collars and arbitrarily small individual scales. An exactly unchanged
normalized complex-hyperbolic exterior outside removed volume v has cap
energy 2v, since H\_CH=2. The classical shrinking-core upper sequence
lies in this class: choose its gluing scales diagonally so that the
defect approximation error is o(h\_j) and cap energy is o(A\_jh\_j). It
then matches coefficient 4 asymptotically as h\_j-\textgreater0. This
proves coefficient sharpness \textbf{within that replacement class}, not
among all metrics on M.

Actual spinor budgets alone were not shown to produce these C2 collars:
they do not control the full opposite Weyl curvature or injectivity
radius. No such geometric conclusion is needed for the weak-Bianchi
proof above.

\hypertarget{retaining-the-signed-current-an-independent-exact-identity}{%
\subsection{7. Retaining the signed current: an independent exact
identity}\label{retaining-the-signed-current-an-independent-exact-identity}}

This records the parallel flux route without presenting it as a second
proof of Theorem 3. Let Gamma=Im zeta and F\_A=iF, with F real and
closed. Use the cutoff chi of Section 1, and define, for rho\textless1,

\[
 B_\chi(\rho)=-\int_\rho^1\frac{\chi'(t)}{t^2}dt,
 \qquad C_\chi(\rho)=-\int_\rho^1\frac{\chi'(t)}t dt.
\]

Both functions are constant below 1/2 and zero above 1. Cauchy--Schwarz
against the positive measure -chi'(t)dt proves C\_chi\^{}2\textless=chi
B\_chi. Their derivatives are chi'/rho\^{}2 and chi'/rho, respectively.
On the closed manifold, (1.6) integrates to int{[}B\_chi dGamma wedge
dGamma-C\_chi dGamma wedge F{]}. Combining with the central curvature in
(1.3) and completing a square gives

\[
\begin{aligned}
16\pi^2c_2(S^+)[M]=\int\bigl[&4B_\chi(|Z^+|^2-|Z^-|^2)
 +\lambda(|F^+|^2-|F^-|^2)\\
 &+2\chi(|\Omega^-|^2-|\Omega^+|^2)\bigr]d\mu+R_{\rm off},
\end{aligned}                                                     \tag{7.1}
\]

where Z=dGamma-C\_chi F/(2B\_chi),
lambda=chi-C\_chi\^{}2/B\_chi\textgreater=0 and R\_off=2int chi' d rho
wedge(b wedge F12-bar b wedge F21). Where B\_chi=0 all weighted terms
are defined by their continuous limit; C\_chi=chi=0 there. The identity
extends through spinor zeros because Gamma is globally smooth and the
weights are constant near rho=0; the off-diagonal expression is
supported away from those zeros.

The central anti-self-dual terms in (7.1) have nonpositive sign. This
does not make the full topological flux nonpositive: its self-dual
current energy remains. The new estimate (4.14), however, now bounds
that current on the beta=2 square-controlled branch. Indeed dGamma is
the sum of a quadratic spinor-gradient term and
Im\textless Phi,F\^{}\{S+\}Phi\textgreater; (4.4)--(4.5), (4.8) and
(4.14) bound its L2 energy by CTheta. Thus the signed identity is
consistent with the linear positive-charge bound; it supplies no
independent zero-end vanishing theorem.

\hypertarget{continuing-the-zero-end-system-torsion-length-and-a-coupled-balance}{%
\subsection{8. Continuing the zero-end system: torsion length and a
coupled
balance}\label{continuing-the-zero-end-system-torsion-length-and-a-coupled-balance}}

This section continues the zero-end problem after the linear theorem.
All norms here are normalized by V. Assume the inherited actual supply
with a\^{}2\textgreater=T\_+, beta\textgreater=4, and
eta=kappa+epsilon\textless=1. Write P=nabla\_A Phi, H=nabla\_A P,
t=t\_Phi and m=\textbar t\textbar\^{}2. Its uniform budget gives

\[
 \|P\|_4+\|H\|_2+\|t\|_4+\|\nabla t\|_2\le C,\quad
 \|\zeta\|_2+\|d\rho\|_2\le C\sqrt\eta,
\] \[
 \|\rho-\rho_\beta\|_2+\|U-U_\beta\|_2
 +\|s+4\|_2+\|\operatorname{Ric}^{\circ}\|_2\le C\sqrt\eta,
\]

where rho\_beta and U\_beta are as in Section2. Set

\[
 m_\beta=-U_\beta\rho_\beta^3/2
 =\frac{16[4(1+1/\beta)]^{3/2}}{81}(1-2/\beta),\quad
 S=m_\beta+2\rho^2|P|^2,\qquad u=t/\sqrt S.            \tag{8.1}
\]

For beta\textgreater=4, m\_beta has uniform positive lower and upper
bounds. All these fields are globally smooth, including at spinor zeros.
The pointwise Dirac-compatible decomposition used in Section3 gives

\[
       0\le d:=2\rho^2|P|^2-m\le4\rho|\zeta|^2.      \tag{8.2}
\]

Indeed write P=alpha v\_1+gamma\_N v\_2+gamma\_L v\_2 using the
orthogonal kernel decomposition there. Then
m=2rho\^{}2\textbar gamma\_N\textbar\^{}2,
\textbar gamma\_L\textbar\textless=\textbar alpha\textbar{} and
\textbar alpha\textbar\^{}2=\textbar zeta\textbar\^{}2/rho. The
inequality extends through zeros. Thus S=m\_beta+m+d,
\textbar u\textbar\textless=1, and avg d\textless=C eta. The inherited
fourth-order and same-solution mass estimates imply

\[
 \fint S=2m_\beta+O(\sqrt\eta),\qquad \fint S^2\le C. \tag{8.3}
\]

\textbf{Proposition 6 (actual torsion-length concentration).} Uniformly
under these assumptions,

\[
 \fint\frac{(S-2m_\beta)^2}{S}\le C\sqrt\eta,\qquad
 \|m-m_\beta\|_{3/2,\mathrm{avg}}\le C\eta^{1/4},
                                                               \tag{8.4}
\] \[
 \frac{\operatorname{Vol}\{m\le m_\beta/4\}}V
                                      \le C\sqrt\eta. \tag{8.5}
\]

In particular, along beta-\textgreater infinity, eta-\textgreater0,
\textbar t\textbar\^{}2 tends to 128/81 strongly in every normalized Lq,
1\textless=q\textless2, not merely in mean.

\textbf{Proof.} Test the actual equation Delta
rho=-2\textbar P\textbar\^{}2-U rho/2 against a=2rho\^{}2/S. Direct
differentiation cancels its leading terms:

\[
 da=\frac{4m_\beta\rho}{S^2}d\rho
       -\frac{8\rho^4}{S^2}\operatorname{Re}\langle P,H\rangle.
\]

The coefficients of d rho and H are uniformly bounded. For the latter,
put z=rho\textbar P\textbar, so
rho\^{}4\textbar P\textbar/S\^{}2=rho\^{}3
z/(m\_beta+2z\^{}2)\^{}2\textless=C. Hence
\textbar\textbar da\textbar\textbar\_2\textless=C and
\textbar avg\textless da,d rho\textgreater\textbar\textless=C sqrt eta.
The equation gives

\[
 \fint(2-2m_\beta/S+U\rho^3/S)=O(\sqrt\eta).
\]

Replace U rho\^{}3 by U\_beta rho\_beta\^{}3=-2m\_beta, at cost O(sqrt
eta). It follows that avg(1/S)=1/(2m\_beta)+O(sqrt eta). The exact
nonnegative Jensen defect avg{[}(S-2m\_beta)\^{}2/S{]}=avg
S-4m\_beta+4m\_beta\^{}2 avg(1/S) and (8.3) prove the first part of
(8.4).

For its stronger norm assertion, test instead with b=rho\^{}2 sqrt S.
Direct differentiation gives
\textbar db\textbar\textless=C{[}(1+\textbar P\textbar)\textbar d
rho\textbar+\textbar H\textbar{]}. Since \textbar d
rho\textbar\textless=4\textbar P\textbar, the inherited L4/Hessian
budget still gives \textbar\textbar db\textbar\textbar\_2\textless=C.
Therefore

\[
 \fint(S-2m_\beta)\sqrt S=O(\sqrt\eta).
\]

The replacement error is at most \textbar\textbar U
rho\^{}3+2m\_beta\textbar\textbar\_2 \textbar\textbar sqrt
S\textbar\textbar\_2\textless=C sqrt eta. Subtract sqrt(2m\_beta) times
the first identity of (8.3). With G=(S-2m\_beta)(sqrt
S-sqrt(2m\_beta))\textgreater=0, obtain avg G\textless=C sqrt eta.
Holder yields

\[
 \fint|S-2m_\beta|^{3/2}
 \le(\fint G)^{3/4}
       [\fint(\sqrt S+\sqrt{2m_\beta})^3]^{1/4}
 \le C\eta^{3/8}.
\]

The last factor is bounded by (8.3). From (8.2), interpolation of the
small L2 and bounded L4 norms of zeta gives
\textbar\textbar d\textbar\textbar\_\{3/2\}\textless=C eta\^{}\{1/3\}.
Since m-m\_beta=S-2m\_beta-d, (8.4) follows. Interpolate with the
bounded L2 norms to obtain every q\textless2, not q=2.

Finally, split \{m\textless=m\_beta/4\} into
\{d\textgreater m\_beta/4\}, of volume O(eta)V, and its complement. On
the complement S\textless=3m\_beta/2 and
(S-2m\_beta)\^{}2/S\textgreater=m\_beta/6; (8.5) follows. QED.

\hypertarget{using-both-the-dirac-and-actual-curvature-equations}{%
\subsubsection{8.1 Using both the Dirac and actual curvature
equations}\label{using-both-the-dirac-and-actual-curvature-equations}}

The equations of Sections2--3 imply an additional quantitative identity:

\[
 \boxed{\fint\left(\rho^3|\nabla u|^2-
            \frac{\rho^3m|dS|^2}{4S^3}\right)
       =\rho_\beta^3(1-1/\beta)+O(\sqrt\eta).}         \tag{8.6}
\]

In particular, the left side tends to 64/27 at the zero endpoint. Here
is the proof, including the critical tests.

Write Dt=R\_D+E, with
\textbar R\_D\textbar\textless=C\textbar zeta\textbar\textbar P\textbar{}
and \textbar\textbar E\textbar\textbar\_2\textless=C sqrt eta. Let w=1/S
and W=rho\^{}3/S. Testing the twisted Lichnerowicz identity against Wt
gives, with real Hermitian pairings,

\[
\begin{aligned}
 \int W|Dt|^2+\int\langle Dt,c(dW)t\rangle
  ={}&\int W|\nabla t|^2+\tfrac12\int\langle dW,dm\rangle\\
    &+\tfrac14\int sWm
      +\tfrac32\int W\langle c(F_A^-)t,t\rangle.       
\end{aligned}\tag{8.7}
\]

The two left integrals are O(sqrt eta V): use
rho\^{}3\textbar P\textbar\^{}2/S\textless=C, and the bounds
\textbar t\textbar\textbar dW\textbar+\textbar P\textbar\textbar t\textbar\textbar dW\textbar\textless=C(\textbar d
rho\textbar+\textbar H\textbar). Their coefficients reduce to bounded
functions of z=rho\textbar P\textbar{} and S=m\_beta+2z\^{}2, so no
division by rho occurs.

The gradient term containing d(rho\^{}3) is also O(sqrt eta V), since
rho\^{}2\textbar t\textbar/S is bounded. The rest completes exactly as

\[
 w|\nabla t|^2+\tfrac12\langle dw,dm\rangle
   =|\nabla(t/\sqrt S)|^2-\frac{m|dS|^2}{4S^3}.       \tag{8.8}
\]

The scalar term is -int rho\^{}3 m/S+O(sqrt eta V). Substitute the
actual curvature equation (3.1) in the last term of (8.7). The identity
c(q\_-(t))t=(m/2)t gives

\[
 -\tfrac32\int m^2/S
 +3\int\frac{\rho^2}{S}\langle c((d\zeta)^-)t,t\rangle
 +\tfrac32\int\frac{\langle c(\mathcal R)t,t\rangle}{S}. \tag{8.9}
\]

The quartic term is negative in the original Lichnerowicz equation. The
last integrand is bounded by
C(\textbar Ric\^{}circ\textbar+\textbar zeta\textbar\textbar P\textbar),
so its normalized integral is O(sqrt eta). For the middle term, Clifford
duality gives
\textless c(xi)t,t\textgreater=4\textless xi,q\_-(t)\textgreater{} for
imaginary anti-self-dual xi. Integrating by parts turns it into

\[
       12\int\left\langle\zeta,
              \delta\left(\frac{\rho^2q_-(t)}S\right)\right\rangle.
\]

The test is smooth globally. Its derivative satisfies

\[
 \left|\nabla\left(\frac{\rho^2q_-(t)}S\right)\right|
       \le C(|\nabla t|+|H|+|d\rho|).                \tag{8.10}
\]

Indeed \textbar q(t)\textbar\textless=Cm, \textbar nabla
q(t)\textbar\textless=C\textbar t\textbar\textbar nabla t\textbar, and
differentiating S reduces the needed coefficients to
rho\^{}2\textbar t\textbar/S, rho m/S,
rho\^{}3m\textbar P\textbar\^{}2/S\^{}2, and
rho\^{}4m\textbar P\textbar/S\^{}2. All are bounded by
m\textless=2rho\^{}2\textbar P\textbar\^{}2 and bounded rho. Thus this
current term is O(sqrt eta V), without assuming the unweighted critical
Dirac residual small.

Combining (8.7)--(8.10) proves

\[
 \fint\left(\rho^3|\nabla u|^2-\frac{\rho^3m|dS|^2}{4S^3}\right)
  =\fint\left(\frac{\rho^3m}{S}+\frac{3m^2}{2S}\right)
      +O(\sqrt\eta).
\]

From S=m\_beta+m+d, avg d=O(eta), and the reciprocal test,
avg(m/S)=1/2+O(sqrt eta) and avg(m\^{}2/S)=m\_beta/2+O(sqrt eta). For
the latter replace m by S-m\_beta at cost at most 2 avg d and use
(S-m\_beta)\^{}2/S=S-2m\_beta+m\_beta\^{}2/S. Replace rho\^{}3 by
rho\_beta\^{}3 in the first potential term, whose multiplier m/S is
bounded. The result is
rho\_beta\^{}3/2+3m\_beta/4=rho\_beta\^{}3(1-1/beta), proving (8.6).

Both derivative integrals have uniform budgets; in particular rho\^{}3
m\textbar dS\textbar\^{}2/S\^{}3\textless=C(\textbar d
rho\textbar\^{}2+\textbar H\textbar\^{}2). Since the subtracted term is
nonnegative, the zero-end consequence is

\[
             \liminf\fint|\nabla u|^2\ge1/27.        \tag{8.11}
\]

We used rho\^{}3\textless=64, not an unjustified replacement of rho by
its limiting constant inside a derivative integral. The more precise
weighted lower bound is liminf avg rho\^{}3\textbar nabla
u\textbar\^{}2\textgreater=64/27.

This is a full coupled identity at the available norms and a strong
subcritical length statement. It is not zero-freeness, critical L2
convergence of m, or a zero-end gap. The derivative defect in (8.6) need
not vanish; curvature/topological charge may still concentrate on the
shrinking low-torsion region. That localized charge, rather than mean
mass alone, remains to be controlled.

There is a concrete chirality distinction at this remaining step. For
B=S\^{}- tensor L the signed density is

\[
 16\pi^2c_2(B)=
 \left(9|F_A^+|^2-9|F_A^-|^2
       +2|W^-|^2+s^2/24-\tfrac12|\operatorname{Ric}^{\circ}|^2\right)d\mu.
                                                               \tag{8.12}
\]

Thus an ideal Ricci-flat core with W\^{}+=F\_A=0 has \emph{positive}
bulk density \textbar W\^{}-\textbar\^{}2/(8pi\^{}2) for this bundle.
The positive-spinor negative-charge argument does not transfer to it.
This is a bulk sign calculation, not a complete marked filling or an
actual counterexample. The next sufficient exclusion would have to
control this positive curvature contribution together with the boundary
flux on the low-torsion region in (8.5), or a construction would have to
realize that entire coupled region and its exterior. Small volume by
itself settles neither alternative.

END COMPLETE SOURCE

\begin{center}\rule{0.5\linewidth}{0.5pt}\end{center}

\hypertarget{source-2-20260927_conformal_separation_and_zero_end.md}{%
\subsection{Source 2:
20260927\_conformal\_separation\_and\_zero\_end.md}\label{source-2-20260927_conformal_separation_and_zero_end.md}}

Original path:
continuation/20260927\_conformal\_separation\_and\_zero\_end.md

SHA-256:
8f37b7ceb8fce5ce2f33ad6a87f5a8afcfe4cfd0b19282ed6ea8ae3689aebca9

BEGIN COMPLETE SOURCE

\hypertarget{conformal-separation-realized-linear-windows-and-the-zero-end-system}{%
\section{Conformal separation, realized linear windows, and the zero-end
system}\label{conformal-separation-realized-linear-windows-and-the-zero-end-system}}

27 September 2026. This is a mathematical continuation of
\href{20260927_mixed_current_and_signed_cores.md}{the linear-cost
proof}, not a new export or a replacement for the inherited manuscripts.
The beta=2 refined Hessian theorem is the analytic premise of Part I.
Part II continues the actual zero-Weyl equations and proves geometric
endpoint exclusions under their stated hypotheses. The unrestricted
integral-error endpoint gap remains open.

Throughout, manifolds are closed, connected, oriented, smooth and
four-dimensional unless a local limit is explicitly specified. Curvature
norms and actual conformal monopole solvability have exactly the
conventions of the
\href{../mathematics/current/line1_functional_pencil.md}{functional
pencil}. In particular, one genuine Spin\^{}c structure is fixed while
the metric varies within its conformal class. A different original
metric may select a different structure. Constants are universal unless
a dependence is displayed.

\hypertarget{part-i.-removing-scalar-variance-and-enlarging-the-actual-energy-window}{%
\subsection{Part I. Removing scalar variance and enlarging the actual
energy
window}\label{part-i.-removing-scalar-variance-and-enlarging-the-actual-energy-window}}

\hypertarget{a-scale-invariant-scalar-slack-theorem}{%
\subsubsection{1. A scale-invariant scalar-slack
theorem}\label{a-scale-invariant-scalar-slack-theorem}}

For one such actual structure, put

\[
\begin{gathered}
 p=(a_g^+)^2>0,\quad T=2\chi+3\tau,\quad
 P=\int|\operatorname{Ric}^{\circ}|^2,\quad S_2=\int s^2,\\
 D=\frac32S_2+18\int|W^+|^2-72\pi^2p,\quad
 \Sigma=S_2-32\pi^2p,\quad Y=(T-a^2)_+,\quad \Pi=P+D.
\end{gathered}                                                     \tag{1.1}
\]

The actual scalar and functional inequalities give
Sigma,D\textgreater=0.

\textbf{Theorem 1.} For every such metric, with no small-error
condition,

\[
 \boxed{\Sigma\le C\bigl(\Pi+\sqrt{\Pi Y}\bigr),\qquad
 (p-T)_+\le C_T\bigl(\Pi+\sqrt{\Pi Y}\bigr).}             \tag{1.2}
\]

In particular, when a\^{}2\textgreater=T,

\[
 \boxed{\Sigma\le C(P+D),\qquad (p-T)_+\le C_T(P+D).}     \tag{1.3}
\]

There is no scalar-variance term, normalized target Ricci error,
diameter, Sobolev constant or noncollapse assumption in (1.2).

\textbf{Proof.} The positive actual projection excludes a nonnegative
constant scalar-curvature representative: the unperturbed spinor
Weitzenboeck formula would force the spinor and the self-dual
determinant curvature to vanish, contradicting p\textgreater0. The
Yamabe representative in {[}g{]} therefore has negative constant scalar
curvature. Rescale it to obtain s\_hat=-4, and denote it by g\_hat. The
same actual structure remains available throughout this conformal class.
Hodge star on two-forms is conformally invariant, so p is unchanged.
Class squares and Y are unchanged as well.

For clarity, the conformal energy comparison is exact, not a claim that
the conformal factor is small. Write g=u\^{}2 g\_hat and use
Delta=d*d.~The scalar transformation law gives s\_g u\^{}3=6 Delta\_hat
u-4u. Squaring and integrating, then using int u\^{}\{-1\}Delta u=-int
u\^{}\{-2\}\textbar du\textbar\^{}2, gives

\[
 \delta:=S_2-\widehat S_2
 =48\int_{\widehat g}u^{-2}|du|^2
       +36\int_{\widehat g}(u^{-1}\Delta u)^2\ge0.        \tag{1.4}
\]

Conformal invariance of the Weyl energy and Chern--Weil give

\[
 P-\widehat P=\delta/12,\quad
 D-\widehat D=3\delta/2,\quad
 \Sigma-\widehat\Sigma=\delta.                          \tag{1.5}
\]

Both hatted energies P\_hat,D\_hat are nonnegative. On g\_hat the
normalized Einstein residual satisfies r\_hat V\_hat=P\_hat. If
(P\_hat+D\_hat)/V\_hat\textless=1, the refined actual theorem and its
square-shortfall version in Sections 4--5 of the preceding proof give

\[
 \widehat\Sigma\le C_0\left(\widehat\Pi+
                             \sqrt{\widehat\Pi Y}\right).
\]

The selectable smoothing allowance is sent to zero for this fixed metric
and actual structure. If instead Pi\_hat\textgreater V\_hat, the
elementary inequality
0\textless=Sigma\_hat=16V\_hat-32pi\^{}2p\textless=16V\_hat gives the
same estimate with constant 16. Thus it holds in both cases, with
C\textgreater=max(C\_0,16). Now Pi=Pi\_hat+(19/12)delta, and (1.5)
immediately transfers the bound to g. Finally use the exact identity

\[
 36\pi^2(p-T)=\frac92P+\frac98\Sigma-D.                 \tag{1.6}
\]

This proves (1.2). QED.

The conformal identities were already available in
\href{../mathematics/current/line1_energy_selection.md}{energy
selection, Section 4} and
\href{../mathematics/current/line1_yamabe_endpoint.md}{Yamabe
normalization, Section 1}. The new inference is their combination with
the refined beta=2 theorem. No pointwise Ricci preservation under a
conformal change is used or claimed. Nor does (1.2) transfer the hatted
spinor's estimates to a particular weighted spinor on the original
metric: its conclusions are curvature energies.

\hypertarget{an-unconditional-lower-ricci-energy-wall}{%
\subsubsection{2. An unconditional lower Ricci-energy
wall}\label{an-unconditional-lower-ricci-energy-wall}}

Assume, for the metric under consideration,

\[
 p\ge A>0,\qquad T=A(1-h),\qquad h>0.
\]

Write

\[
 R=\int|\operatorname{Ric}|^2,\quad
 e=\frac{R}{8\pi^2A},\quad
 \mathcal J=\frac{4P}{R},\quad z=p/A-1\ge0,\quad y=Y/A.
\]

The scalar budget makes R positive. Always 0\textless=J\textless=4 and
J\textless=E, where E is the project's relative Einstein defect. The two
exact reductions are

\[
 \frac DA=36\pi^2[e-1-h-2z],\qquad
 \frac PA=2\pi^2e\mathcal J.                           \tag{2.1}
\]

In particular e\textgreater=1+h+2z. Theorem 1 now has the following
consequence.

\textbf{Theorem 2.} If a\^{}2\textgreater=T, there is c\_0\textgreater0,
independent of h and e, such that

\[
 \boxed{e(1+\mathcal J/18)
       \ge1+(1+c_0)h+(2+c_0)z.}                       \tag{2.2}
\]

In general there is c\_s\textgreater0 such that

\[
 \boxed{e(1+\mathcal J/18)
       \ge1+h+2z+c_s\frac{(h+z)^2}{h+z+y}.}            \tag{2.3}
\]

The estimates require no upper energy budget from the other chirality.
In particular every non-Ricci-flat Einstein candidate in the
square-controlled case satisfies e\textgreater=1+(1+c\_0)h.

\textbf{Proof.} In the square-controlled case, (1.3) gives
A(h+z)\textless=C\_T(P+D). Substitute (2.1), and take c\_0=(36pi\^{}2
C\_T)\^{}\{-1\}. For the general case put x=Pi/A and v=h+z. The
inequality v\textless=C\_T(x+sqrt(xy)) implies x\textgreater=c
v\^{}2/(v+y): if x\textgreater=v this is immediate; otherwise square and
use x\textless v. Substitution proves (2.3). QED.

For 0\textless h\textless=1/3, (2.2) also gives J+(e-1-h)\textgreater=c
h. Indeed if e-1-h\textgreater=h/2 there is nothing to prove; otherwise
e\textless=3/2 and e-1-h+J/12\textgreater=c\_0h. This is a global
energy-versus-anisotropy statement, not merely an exact Einstein
inequality.

\hypertarget{a-stronger-chiral-window-and-the-small-h-variational-problem}{%
\subsubsection{3. A stronger chiral window and the small-h variational
problem}\label{a-stronger-chiral-window-and-the-small-h-variational-problem}}

Put U=3 min(T\_+,T\_-)/A. Chern--Weil yields directly

\[
 e(1-\mathcal J)\le U.                                 \tag{3.1}
\]

Unlike the older formula with E, no scalar variance has been discarded
here. Let L\_h=1+(1+c\_0)h. If U\textgreater0, combining (2.2) and (3.1)
gives

\[
 \boxed{\mathcal J\ge
       \frac{(L_h-U)_+}{L_h+U/18}.}                    \tag{3.2}
\]

For J\textless1 this follows by eliminating e. For J\textgreater=1 the
right side is at most one, so it remains true. If U\textless=0, (3.1)
itself gives J\textgreater=1. Thus for 0\textless h\textless=1/3,

\[
 \mathcal J+(U-1-h)_+\ge c h.                          \tag{3.3}
\]

The analogous square-shortfall result has right side c h\^{}2/(h+y), or
c h/(1+L) when y\textless=Lh. All statements also hold with E in place
of J.

On the already constructed actual small-h family in
\href{20260926_boundary_compatibility.md}{Section 3 of the boundary
continuation}, U\_j=1+h\_j+O(h\_j\^{}2), y\_j=0. Consequently

\[
 \boxed{c_1h_j\le\inf_g\mathcal J(g)
                  \le\frac{4h_j}{1+h_j}.}              \tag{3.4}
\]

The upper bound uses the previously justified classical shrinking-core
metrics and J\textless=E. The lower bound applies to every smooth
metric. The linear order is therefore forced by traceless Ricci itself;
scalar variance cannot be the sole payment of the topological cost.

Here is a useful exact variational comparison. From (1.5),

\[
 R-\widehat R=\delta/3,\qquad
 \mathcal J(g)=
 \frac{\mathcal J(\widehat g)\widehat R+\delta/3}
      {\widehat R+\delta/3}.                           \tag{3.5}
\]

Thus J(g) lies between J(g\_hat) and 1. Since the normalized negative
Yamabe metric has constant scalar curvature, E(g\_hat)=J(g\_hat).
Whenever J(g)\textless1,

\[
 \mathcal E(\widehat g)=\mathcal J(\widehat g)
                 \le\mathcal J(g)\le\mathcal E(g).      \tag{3.6}
\]

If the manifold has the stipulated actual supply for every metric and
inf\_g J(g)\textless1, it follows that

\[
                    \inf_g\mathcal E(g)=\inf_g\mathcal J(g).
                                                               \tag{3.7}
\]

This applies to the small-h family for all sufficiently large j. It
concerns infima, not attainment, pointwise closeness, or uniqueness of
Einstein metrics.

\hypertarget{conformal-minimization-and-minimizing-sequences}{%
\paragraph{3.1 Conformal minimization and minimizing
sequences}\label{conformal-minimization-and-minimizing-sequences}}

The comparison can be sharpened without further analytic input. Fix a
negative conformal class, and put J\_Y=J(g\_hat), R\_Y=R(g\_hat).
Formula (3.5) shows

\[
             \inf_{g\in[g_0]}\mathcal J(g)=\min\{J_Y,1\}.
                                                               \tag{3.8}
\]

To justify the possibly unattained endpoint 1, choose a coordinate ball
and a compactly supported smooth cutoff chi which is one on a smaller
ball. For fixed sufficiently small epsilon\textgreater0 let
u\_n=1+epsilon chi sin(nx\^{}1). These are uniformly positive, bounded
conformal factors. The n\^{}2 term in Delta\_hat u\_n has nonzero
squared integral of order n\^{}4 on the smaller ball, while the
lower-order terms are O(n); hence the nonnegative expression (1.4) tends
to infinity. Formula (3.5) tends to 1. If J\_Y\textless1, (3.8) is
attained exactly by constant conformal rescalings of the Yamabe metric.
If J\_Y=1, every metric in the class has J=1. If J\_Y\textgreater1, the
infimum 1 is not attained.

If every conformal class on M is negative (as follows from the
all-metric actual positive supply), write alpha=inf\_g E(g). Since
J\textless=E, while (3.6) would produce an E-value below alpha from any
J-value below min(alpha,1),

\[
                 \boxed{\inf_g\mathcal J(g)=\min\{1,\alpha\}.}
                                                               \tag{3.9}
\]

When alpha\textless1 there is an exact useful excess decomposition.
Define L\_s=int(s-s\_bar)\^{}2 and delta=S\_2-S\_Y. Since E=(4P+L\_s)/R
and E(g\_hat)=J\_Y\textgreater=alpha, (3.5) gives

\[
 \boxed{\mathcal E(g)-\alpha=
   \frac{R_Y}{R}(J_Y-\alpha)
   +(1-\alpha)\frac{\delta}{3R}+\frac{L_s}{R}.}          \tag{3.10}
\]

Every summand is nonnegative. In particular
delta/R\textless=3(E-alpha)/(1-alpha) and L\_s/R\textless=E-alpha. Every
E-minimizing sequence therefore has delta/R-\textgreater0,
R\_Y/R-\textgreater1, and J\_Y-\textgreater alpha. Its late terms may be
replaced by their scalar-normalized Yamabe metrics with no larger E.
This reduces the low-defect variational problem to constant-scalar
metrics and controls the conformal energy discarded in the reduction. It
does not supply convergence of conformal factors, geometric compactness,
or attainment of the infimum.

\hypertarget{the-remaining-sharp-coefficient-is-a-specific-remainder-problem}{%
\paragraph{3.2 The remaining sharp coefficient is a specific remainder
problem}\label{the-remaining-sharp-coefficient-is-a-specific-remainder-problem}}

For the realized family with U=1+h+O(h\^{}2), h-\textgreater0, the upper
bound tends to zero. The preceding reduction therefore permits
restricting the variational infimum to the actual scalar -4 Yamabe
representatives. Put

\[
 x=\frac{P}{8\pi^2A},\quad d=\frac{D}{36\pi^2A},\quad
 \sigma=\frac{\Sigma}{32\pi^2A},\quad z=p/A-1.
\]

All are nonnegative, and exact identities give

\[
\begin{gathered}
 e=1+h+2z+d,\qquad V/(2\pi^2 A)=1+z+\sigma,\\
 x=h+z+d-\sigma,\qquad
 \mathcal E=\mathcal J=\frac{4(h+z+d-\sigma)}{1+h+2z+d},\\
                 d+2z\le4x+(U-1-h).
\end{gathered}                                                     \tag{3.11}
\]

The minimization is over triples realized by these actual metrics and
classes, not arbitrary nonnegative numbers satisfying (3.11).

For the same class take an actual ordinary solution with weight one.
This is available for the present actual family. It is a separately
supplied solution, not the variable-weight spinor silently rescaled.
Write rho=\textbar Phi\textbar\^{}2, K=int\textbar nabla\_A
Phi\textbar\^{}2, Q=int(4-rho)\^{}2, and
J\_+=\textbar\textbar F\_A\^{}+-(F\_A\^{}+)\_H\textbar\textbar\_2\^{}2.
The maximum principle gives rho\textless=4, and the same-solution
identities give

\[
 4K=\int(4-\rho)\rho,\quad
 \int\rho^2=32\pi^2p+8J_+,\quad
 \boxed{\Sigma=8K+Q+8J_+.}                             \tag{3.12}
\]

The last formula is the ordinary-weight specialization of the
\href{20260926_actual_topological_cost.md}{actual remainder identity},
not a new gauge estimate. It holds for every such solution, so changing
the solution does not alter the sum of these three losses.

Along any actual candidate sequence with E=O(h), (3.1) and
e\textgreater=1+h+2z+d imply e=1+O(h). Thus x,z,d,sigma=O(h), and (3.11)
gives E/h=4-4(sigma-z-d)/h+O(h). Consequently

\[
 \boxed{\liminf\frac{\mathcal E}{h}\ge4
 \quad\Longleftrightarrow\quad
 \frac{[K+Q/8+J_+-D/9-4\pi^2(p-A)]_+}{Ah}\longrightarrow0.}
                                                               \tag{3.13}
\]

Equivalently, the still-unproved first-order estimate is

\[
                 \Sigma\le\frac89D+32\pi^2(p-A)+o(Ah)
                                                               \tag{3.14}
\]

on these candidate sequences. Obtaining it for all such sequences would
determine coefficient4, in view of the known upper construction. It is
not asserted for arbitrary manifolds or as a consequence of Theorem 1.
On a sharp minimizing sequence it asks for sigma=z+d+o(h), not for each
slack separately to vanish. Controlling K alone loses the other two
nonnegative terms in (3.12).

There is an actual obstruction to reaching this coefficient merely by
optimizing C\_T in (1.3). Take a product of closed surfaces of genera
g\_1,g\_2\textgreater=2, with Gaussian curvatures -(1+t), -(1-t),
0\textless t\textless1. Set A\_0=T=8(g\_1-1)(g\_2-1). The product is a
negative Yamabe cscK metric with s=-4 and actual canonical input. Direct
product calculations give

\[
 p=\frac{A_0}{1-t^2},\quad
 P=\frac{8\pi^2 A_0t^2}{1-t^2},\qquad D=\Sigma=0.
                                                               \tag{3.15}
\]

The canonical ordinary solution is parallel with rho=4, so K=Q=J\_+=0.
Thus p-T=P/(8pi\^{}2), forcing C\_T\textgreater=1/(8pi\^{}2) and
c\_0\textless=2/9. The slope certified by (3.2) when U=1+h+O(h\^{}2) is
18c\_0/19, at most 4/19. Improving that isotropic constant alone cannot
yield 4. This product test has h=0 for the common budget A=A\_0; it is
not a counterexample in the tight small-h window or to (3.14).

Finally Yamabe replacement does not improve the known classical upper
coefficient. For each fixed surgery manifold it decreases Ricci energy,
while LeBrun's absolute lower bound prevents it falling below the same
infimum. Hence e-\textgreater1+h, and (2.1) forces z,d-\textgreater0.
The classical completed square bound
Sigma\textless=4D/9+(4sqrt(2)/3)sqrt(RD) then gives Sigma-\textgreater0.
A sufficiently accurate diagonal sequence has z,d,sigma=o(h), still
yielding E/h-\textgreater4. The unresolved improvement must therefore
control the actual combined remainder in (3.13), not merely discard
scalar variance.

\hypertarget{a-realized-bf-family-outside-the-old-quadratic-strip}{%
\subsubsection{4. A realized BF family outside the old quadratic
strip}\label{a-realized-bf-family-outside-the-old-quadratic-strip}}

This application changes the region actually excluded by the project. It
is not a claim that these smooth manifolds have never been obstructed
elsewhere. The actual class supply and symplectic geography are
classical; their precise versions are recorded in
\href{../mathematics/current/line1_bauer_furuta_geography.md}{the
existing BF geography proof}. See also
\href{https://arxiv.org/html/math/0111228}{Ishida--LeBrun} and
\href{https://arxiv.org/pdf/math/0306021}{Braungardt--Kotschick, Theorem
4}.

Fix theta=min(c\_0/2,1/2)\textgreater0 and eta\_j=1/j. For sufficiently
large j choose an even integer N\_j, increasing as fast as necessary,
with N\_j\textgreater=j\^{}6 and
N\_j\textgreater=2c(eta\_j\^{}2/2)/eta\_j\^{}2, where c(epsilon) is the
geography constant. Set

\[
 d_j=\lceil\eta_j^2N_j\rceil,\quad a_j=9N_j-d_j,
 \quad A_j=2a_j.
\]

The geography theorem supplies a simply connected minimal symplectic
X\_j with chi\_h(X\_j)=N\_j and c\_1\^{}2(X\_j)=a\_j: indeed
a\_j\textless=(9-eta\_j\^{}2/2)N\_j-c(eta\_j\^{}2/2). Choose N\_j
additionally large enough for the theorem's lattice-range condition. Its
invariants are

\[
 \chi(X_j)=3N_j+d_j,\quad \tau(X_j)=N_j-d_j,\quad
 b_+(X_j)=2N_j-1\equiv3\pmod4,\quad b_-(X_j)=N_j+d_j-1.
\]

For the pair of building blocks define

\[
 u_j=3-18\frac{2\tau(X_j)}{A_j}
     =1+\frac{16d_j}{a_j}=1+O(\eta_j^2),\qquad
 U_j^*=1+(1+\theta)\eta_j,
\] \[
 x_j=\frac{U_j^*-u_j+3\eta_j}{18},\qquad
 y_j=\frac{\eta_j-x_j}{4}.
\]

They are positive for large j, with x\_j/eta\_j-\textgreater(4+theta)/18
and y\_j/eta\_j-\textgreater(14-theta)/72. Take even surgery counts

\[
 k_j=2\lfloor A_jx_j/2\rfloor,\qquad
 m_j=2\lfloor A_jy_j/2\rfloor-4,
\]

and form the actual smooth manifold

\[
 M_j=X_j\#X_j\#k_j\overline{\mathbb {CP}}^{\,2}
                         \#m_j(S^1\times S^3).         \tag{4.1}
\]

\textbf{Theorem 3.} On this family, for all sufficiently large j,

\[
 \inf_g\mathcal J(g)\ge c_2h_j>0,\qquad
 U_j-F(h_j)\sim\theta h_j,
 \qquad h_j\longrightarrow0,                           \tag{4.2}
\]

where h\_j=(4+k\_j+4m\_j)/A\_j and F(h)=(4-sqrt(1-3h))\^{}2/9. The
manifolds have no Einstein metric of any sign, ordinary SW invariants
vanish, and neither orientation admits an almost-complex structure. Both
strict Hitchin--Thorpe inequalities hold.

\textbf{Proof.} Write k\_j=A\_jx\_j-alpha\_j and m\_j=A\_jy\_j-beta\_j-4
with alpha\_j,beta\_j in {[}0,2). The exact topology gives

\[
 h_j=\eta_j-\frac{12+\alpha_j+4\beta_j}{A_j},\qquad
 U_j=U_j^*+\frac{36-15\alpha_j+12\beta_j}{A_j}.          \tag{4.3}
\]

Here the signature is positive for large j, so the smaller chirality
really is T\_-. In particular T\_+\textgreater0 and T\_-\textgreater0.
The correction in U is positive and O(A\_j\^{}\{-1\}); the correction in
h is negative and O(A\_j\^{}\{-1\}). Since
A\_j\^{}\{-1\}=o(eta\_j\^{}2), and F(h)=1+h+h\^{}2+O(h\^{}3),

\[
 U_j-F(h_j)=\theta\eta_j+O(\eta_j^2),\qquad
 L_{h_j}-U_j=(c_0-\theta)\eta_j+o(\eta_j).              \tag{4.4}
\]

BF nonvanishing supplies the actual canonical-sum and exceptional-sign
structures. For the four-input formulation of Ishida--LeBrun's theorem,
one may add auxiliary X\_3=X\_4=K3: the total b\_+ is 4N\_j+4, congruent
to 4 modulo 8, and the theorem's m=2 conclusion concerns just the two
target summands. The auxiliary K3 surfaces are not included in M\_j. As
in the existing proof, averaging the independent exceptional signs
selects, for every metric, an actual class with p\textgreater=A\_j and
a\^{}2=A\_j-k\_j\textgreater=T\_+=A\_j-4-k\_j-4m\_j. Thus (3.2) applies
and proves the first claim in (4.2). Nonzero Einstein constants have
J=0; Ricci-flatness is also excluded directly by the positive actual
scalar budget.

Finally, b\_1(M\_j)=m\_j, b\_+(M\_j)=4N\_j-2, and
b\_-(M\_j)=2N\_j+2d\_j-2+k\_j are all even. An almost-complex
four-manifold has (chi+tau)/4=(1-b\_1+b\_+)/2 integral, which is
impossible here, in either orientation. Ordinary SW vanishing follows
from the positive-b\_+ connected-sum decomposition. Increasing N\_j
ensures distinct b\_+ and hence pairwise nonhomeomorphism. QED.

This is outside the old \emph{quantitative} strip:
(U\_j-F(h\_j))/h\_j\^{}2-\textgreater infinity, so the old inequality
E+(U-F(h))\_+\textgreater=c\_old h\^{}2 does not by itself give a
positive E on this family. The classical optimized energy interval also
allows e between F(h\_j) and U\_j. The new linear wall excludes it. The
old near-1/3 estimate does not apply effectively as h\_j-\textgreater0,
and the old upper wall concerned h\textgreater1/4. No assertion of an
upper near-Einstein construction on these symplectic BF manifolds is
made.

\hypertarget{cover-compatible-inputs-are-retained-not-strengthened}{%
\subsubsection{5. Cover-compatible inputs are retained, not
strengthened}\label{cover-compatible-inputs-are-retained-not-strengthened}}

The preceding arguments extend to the invariant-weight version of the
actual input in
\href{../mathematics/current/round14_twisted_rigidity.md}{the Pin-minus
proof}. Suppose pi:M\_tilde-\textgreater M is a finite regular oriented
cover, g\_tilde=pi\^{}*g, and one ordinary structure upstairs supplies
every positive deck-invariant weight required by the proof. The beta=2
refined argument only uses these weights and therefore holds upstairs
without requiring all other weights.

The normalized negative Yamabe representative upstairs is deck-invariant
by uniqueness of the scalar -4 representative in a negative conformal
class. It descends to M. Every further auxiliary weight used on that
representative is still deck-invariant, so the same actual supply
remains valid. Apply Theorems 1--2 upstairs, divide the extensive
energies and topology by the covering degree, and keep the normalized
quantities e,J,E unchanged. The square-shortfall term divides correctly:
sqrt(Pi\_tilde Y\_tilde)/degree=sqrt(Pi\_per\_sheet Y\_per\_sheet).
There is no degree-dependent loss.

For a genuine Pin-minus double-cover input, the inherited lifting
theorem gives p\_tilde=2p, a\_tilde\^{}2=2a\^{}2 and T\_tilde=2T, so all
statements above apply to the twisted budgets downstairs. This is the
actual lifting input of
\href{https://arxiv.org/pdf/1505.05948}{Ishida--Matsuo--Nakamura}, not a
claim that an arbitrary formal local-system class has solutions. The
conclusion does not require the covering manifold to solve arbitrary
non-invariant weights or metrics.

\hypertarget{part-ii.-the-remaining-zero-weyl-endpoint}{%
\subsection{Part II. The remaining zero-Weyl
endpoint}\label{part-ii.-the-remaining-zero-weyl-endpoint}}

The following continuation uses the actual torsion system and length
estimates in Section 8 of the preceding proof. Unlike Part I, it has not
yet produced a universal zero-end gap. Further results are recorded
below as their proofs close, without turning them into an export or a
completion claim.

\hypertarget{saturation-of-determinant-curvature-not-just-torsion-length}{%
\subsubsection{6. Saturation of determinant curvature, not just torsion
length}\label{saturation-of-determinant-curvature-not-just-torsion-length}}

Use the actual fields and notation of the preceding proof, Section 8:
P\_spin=nabla\_A Phi, rho=\textbar Phi\textbar\^{}2,
zeta=\textless Phi,P\_spin\textgreater, t=t\_Phi,
m=\textbar t\textbar\^{}2, eta=kappa+epsilon, beta\textgreater=4, and
a\^{}2\textgreater=T\_+. Set rho\_beta=2sqrt(4(1+1/beta))/3,
U\_beta=(4/3)(2/beta-1), m\_beta=-U\_beta rho\_beta\^{}3/2, and
S=m\_beta+2rho\^{}2\textbar P\_spin\textbar\^{}2. All norms in this
section are normalized by V. The quadratic map q\_- obeys
\textbar q\_-(t)\textbar\^{}2=m\^{}2/8. Define bounded smooth fields

\[
 Q=\frac{\rho^3q_-(t)}S,\qquad
 A_\beta=\frac{4m_\beta}{\rho_\beta^6}Q.                \tag{6.1}
\]

There is no division by rho or m. S has a universal positive lower
bound.

\textbf{Theorem 4.} Under these actual hypotheses,

\[
 \boxed{\|F_A^-+A_\beta\|_{2,\mathrm{avg}}^2
       +\frac{4\pi^2(a^2-T_+)}V
       \le\frac8{3\beta}(1-2/\beta)+C\sqrt\eta.}       \tag{6.2}
\]

Consequently, along beta-\textgreater infinity and eta-\textgreater0,

\[
 \fint|F_A^-|^2\longrightarrow\frac29,
 \qquad (a^2-T_+)/V\longrightarrow0.                   \tag{6.3}
\]

More quantitatively, for any measurable B subset M,

\[
 \fint_B|F_A^-|^2\le
 C\left(\frac{\operatorname{Vol}(B)}V+\beta^{-1}
                                      +\sqrt\eta\right),        \tag{6.4}
\]

where fint\_B in (6.4) means V\^{}\{-1\}int\_B, not normalization by
Vol(B). Thus determinant ASD energy cannot concentrate at the normalized
volume scale in the hypothetical endpoint sequence.

\textbf{Proof.} Pair the actual curvature equation

\[
 \rho^3F_A^-+2q_-(t)=2\rho^2(d\zeta)^-+\mathcal R
\]

with q\_-(t)/S. The derivative test rho\^{}2 q\_-(t)/S has bounded
normalized L2 derivative by (8.10) of the preceding proof. Integration
by parts against zeta therefore costs O(sqrt eta), and the remainder
costs the same. Hence

\[
 \fint\langle F_A^-,Q\rangle
     =-\fint\frac{m^2}{4S}+O(\sqrt\eta)
     =-m_\beta/8+O(\sqrt\eta).                         \tag{6.5}
\]

The reciprocal scalar test also implies

\[
 \fint |Q|^2=\rho_\beta^6/32+O(\sqrt\eta).             \tag{6.6}
\]

For the error rate, write S=m\_beta+m+d, avg d=O(eta). The Jensen defect
avg{[}(S-2m\_beta)\^{}2/S{]}=O(sqrt eta), together with
S\textgreater=m\_beta, controls avg(1/S-1/(2m\_beta))\^{}2 by O(sqrt
eta). It follows that avg(m/S)\^{}2=1/4+O(sqrt eta): here one also uses
the independently proved mean identity avg(1/S)=1/(2m\_beta)+O(sqrt
eta), and avg(d/S)=O(eta). The squared fluctuation alone would not give
this error rate. Replacing rho\^{}6 by rho\_beta\^{}6 is legitimate
because m/S\textless=1 and rho is bounded with small L2 amplitude error.
This proves (6.6), not merely convergence without a rate.

The actual positive-curvature fit and Chern--Weil give

\[
 \fint|F_A^+|^2=\frac89(1+1/\beta)^2+O(\sqrt\eta),
\] \[
 \frac{4\pi^2T_+}V=\frac23+\frac{16}{3\beta^2}
                                         +O(\sqrt\eta).
\]

Using
int(\textbar F\_A\^{}+\textbar\^{}2-\textbar F\_A\^{}-\textbar\^{}2)=4pi\^{}2a\^{}2
yields the exact upper-budget expansion

\[
 \fint|F_A^-|^2+\frac{4\pi^2(a^2-T_+)}V
 =\frac{2+16/\beta-40/\beta^2}{9}+O(\sqrt\eta).          \tag{6.7}
\]

Completing the square with (6.5)--(6.6) subtracts
m\_beta\^{}2/(2rho\_beta\^{}6)=2(1-2/beta)\^{}2/9. The difference is
exactly 8(1-2/beta)/(3beta), proving (6.2). The square term and the
class-square term are both nonnegative. Since A\_beta is uniformly
bounded, (6.4) follows by the elementary squared triangle inequality.
Its total norm and (6.2) give (6.3). QED.

For a sequence with bounded total volume, the nonnegative integer
(a\^{}2-T\_+)/4=c\_2(S\^{}+){[}M{]} must eventually be zero. This does
not make the particular spinor section zero-free. With unbounded volume,
(6.2) gives a density conclusion only; individual bubbles can still
carry finite energy.

\hypertarget{opposite-weyl-bulk-cancels-leaving-compulsory-signed-transition-flux}{%
\subsubsection{7. Opposite-Weyl bulk cancels, leaving compulsory signed
transition
flux}\label{opposite-weyl-bulk-cancels-leaving-compulsory-signed-transition-flux}}

Let E\_-=S\^{}- tensor L, whose determinant connection is 3A. Its signed
density is

\[
 16\pi^2c_2(E_-)=\bigl(2|W^-|^2+k\bigr)d\mu,
\quad k=9(|F_A^+|^2-|F_A^-|^2)+s^2/24
                                      -|\operatorname{Ric}^{\circ}|^2/2.
                                                               \tag{7.1}
\]

Choose a smooth theta(m) equal to zero for m\textless=m\_beta/16 and one
for m\textgreater=m\_beta/8, with uniformly bounded derivative, and put
chi=1-theta\^{}2. Let
T\_beta=\{m\_beta/16\textless m\textless m\_beta/8\}. The actual
torsion-length theorem gives Vol\{chi!=0\}/V\textless=C sqrt eta. By
Theorem 4 and the other bounded-model fits,

\[
 \fint k=20/3+O(\beta^{-1}+\sqrt\eta),\qquad
 V^{-1}\int\chi|k|\le C(\beta^{-1}+\sqrt\eta).          \tag{7.2}
\]

On \{t!=0\}, let T\_t be the rank-two Chern transgression for
t/\textbar t\textbar, using the same normalization as Section 1 of the
preceding proof. Since chi=1 near every zero, its localization is
globally legitimate:

\[
 \boxed{2\int\theta^2|W^-|^2+\int\theta^2 k
                  =16\pi^2\int d\chi\wedge T_t.}       \tag{7.3}
\]

In particular, even an arbitrarily large positive W\^{}- bulk charge in
the small-torsion region is not an independent payment of the
obstruction: the identical term occurs in the global characteristic
number and cancels from (7.3). The right side must supply a positive
signed flux of order V.

\textbf{Theorem 5.} With xi=\textless t,nabla t\textgreater, the actual
endpoint fields satisfy

\[
 \boxed{C V^{-1}\int_{T_\beta}|\xi|^2|\nabla t|^2
    \ge\frac{20}{3}-C(\beta^{-1}+\sqrt\eta)
                          +V^{-1}\int\theta^2|W^-|^2.}          \tag{7.4}
\]

Thus a hypothetical zero-end sequence necessarily has
V\^{}\{-1\}int\_\{T\_beta\}\textbar\textless t,nabla
t\textgreater\textbar\^{}2\textbar nabla
t\textbar\^{}2\textgreater=c\textgreater0 on regions of vanishing
relative volume. No upper bound for W\^{}- is required for this
statement.

\textbf{Proof.} In the unit t-frame write the connection as
{[}{[}alpha,-bar b{]},{[}b,d{]}{]}. On T\_beta its amplitude is bounded
above and away from zero, so
\textbar alpha\textbar\textless=C\textbar xi\textbar,
\textbar b\textbar\textless=C\textbar nabla t\textbar, and
\textbar dm\textbar\textless=2\textbar xi\textbar. The transgression is

\[
 4\pi^2T_t=-\alpha\wedge F_{22}
   +\tfrac12(b\wedge F_{12}-\bar b\wedge F_{21})
   +\alpha\wedge\bar b\wedge b.
\]

Moreover \textbar chi'\textbar\textless=C theta. The curvature of E\_-
is bounded algebraically by
C(\textbar W\^{}-\textbar+\textbar s\textbar+\textbar Ric\^{}circ\textbar+\textbar F\_A\textbar).
Young's inequality therefore gives, for any fixed
epsilon\_1\textgreater0,

\[
\begin{aligned}
16\pi^2\left|\int d\chi\wedge T_t\right|
 \le{}&\epsilon_1\int\theta^2|W^-|^2
       +C_{\epsilon_1}\int_{T_\beta}|\xi|^2|\nabla t|^2\\
      &+C_{\epsilon_1}\int_{T_\beta}
              (s^2+|\operatorname{Ric}^{\circ}|^2+|F_A|^2).
\end{aligned}                                                     \tag{7.5}
\]

For example the W\^{}- term is bounded by C theta
\textbar W\^{}-\textbar\textbar xi\textbar\textbar nabla t\textbar; the
cubic connection term by C theta \textbar xi\textbar\^{}2\textbar nabla
t\textbar\^{}2. All other curvature terms use the same Young inequality
without needing an absorption sign. The last line is
O((beta\^{}\{-1\}+sqrt eta)V), by (6.4), the positive-curvature and
scalar fits, and the small volume of T\_beta. Take epsilon\_1=1 in
(7.5), use (7.2)--(7.3), and rearrange. QED.

This also forces int\_\{T\_beta\}\textbar nabla
t\textbar\^{}4\textgreater=cV and sup\_\{T\_beta\}\textbar nabla
t\textbar\textgreater=c eta\^{}\{-1/8\} along a zero-end sequence. It
does not prove the needed mixed energy is uniformly integrable. Uniform
integrability of this particular transition density would suffice to
exclude the endpoint; it is not an added premise in any earlier theorem.

\hypertarget{a-local-compatibility-gap-and-collapse-in-measure}{%
\subsubsection{8. A local compatibility gap and collapse in
measure}\label{a-local-compatibility-gap-and-collapse-in-measure}}

This section gives an independent geometric consequence with a
bounded-Ricci hypothesis. It does not replace the no-compactness,
integral-Ricci endpoint question. The two classical inputs are the local
harmonic-radius estimate of
\href{https://arxiv.org/pdf/1406.6534}{Cheeger--Naber, Theorem 1.13 and
Definition 2.9}, and the local nonexistence of negative Einstein ASD
almost-Kaehler metrics in
\href{https://johnarmstrongmaths.com/publications/resources/thesis.pdf}{Armstrong,
Theorem 3.3.5}. The proof below connects them to the actual
approximate-form supply.

Fix K,L,v\textgreater0. Let g be a complete oriented four-metric with
\textbar Ric\textbar\textless=K, let omega be a smooth self-dual
two-form with \textbar omega\textbar\textless=L, and let the constant
parameter t belong to {[}0,1/4{]}. Define

\[
 Z_t=|\operatorname{Ric}+g|^2+
 |W^+-t(|\omega|^2I/3-\omega\otimes\omega)|^2
                  +( |\omega|^2-2)^2+|d\omega|^2.       \tag{8.1}
\]

\textbf{Theorem 6 (local coercivity).} There is
epsilon(K,L,v)\textgreater0 such that

\[
       \operatorname{Vol}B_1(x)\ge v
            \quad\Longrightarrow\quad
       \int_{B_2(x)}Z_t\ge\epsilon(K,L,v).             \tag{8.2}
\]

\textbf{Proof.} Otherwise take a sequence whose local integral tends to
zero. Volume comparison gives a fixed lower volume for B\_\{1/2\}(x).
The local harmonic-radius estimate supplies a point in B\_\{1/2\}(x)
with harmonic radius at least r\_0(K,v)\textgreater0. In a fixed-size
chart compactly contained in B\_1(x), the metrics have uniform
C\^{}\{1,alpha\} and W\^{}\{2,p\} bounds for each finite p. Pass on
smaller charts to C\^{}\{1,alpha'\} and weak W\^{}\{2,p\} convergence.

The local L2 Ricci residual tends to zero. Its pointwise norm is bounded
by K+2, so interpolation gives convergence to zero in every finite Lp.
The harmonic-coordinate metric equation shows the limit is a weak
Einstein metric with Ric=-g, hence smooth. Subtracting that equation
from the sequence and using interior elliptic estimates improves the
metric convergence to strong W\^{}\{2,p\} on smaller charts. In
particular the Weyl operators converge strongly in L2. No Ricci
derivative bound is assumed.

Self-duality gives \textbar delta\_g omega\textbar=\textbar d
omega\textbar. Interior estimates for d+delta in the uniformly
C1-controlled charts give a uniform H1 bound for omega; Rellich
compactness gives strong L2 convergence on a smaller ball. The limit is
weakly closed and coclosed, hence smooth, self-dual, and of squared
length two. Bounded amplitude also makes its polynomial model converge
strongly in L2. Passing to a subsequence t-\textgreater t\_infinity
yields

\[
 \operatorname{Ric}=-g,\quad d\omega=0,\quad |\omega|^2=2,
 \quad W^+=t_\infty(2I/3-\omega\otimes\omega).          \tag{8.3}
\]

If t\_infinity=0 this is the locally forbidden ASD Einstein
almost-Kaehler system. If t\_infinity\textgreater0, the Einstein Bianchi
identity gives div W\^{}+=0. Differentiating (8.3), using constant
length and delta omega=0, gives omega\_\{ij\}nabla\^{}i omega\_\{kl\}=0.
The two-form is nondegenerate, so nabla omega=0. The two-form
Weitzenboeck equation now implies W\^{}+omega=(s/6)omega=-(2/3)omega.
But (8.3) gives W\^{}+omega=-(4t\_infinity/3)omega, forcing
t\_infinity=1/2, outside {[}0,1/4{]}. Both cases contradict existence of
the limit. QED.

\textbf{Corollary 7 (collapse in measure).} On a closed manifold put
S\_v=\{x:Vol B\_1(x)\textgreater=v\}. Then

\[
 \boxed{\frac{\operatorname{Vol}(S_v)}V
       \le C(K,L,v)\fint Z_t.}                        \tag{8.4}
\]

Indeed integrate (8.2) over S\_v. Fubini and the Bishop upper bound
b\_K(2) for a radius-two ball give

\[
 \epsilon\operatorname{Vol}(S_v)
 \le\int_M Z_t(y)\operatorname{Vol}(S_v\cap B_2(y))d\mu(y)
 \le b_K(2)\int_M Z_t.
\]

There is no packing multiplicity, total-volume upper bound or diameter
assumption. In particular, for the actual supplier, beta\textgreater=4,
p\textgreater0, and \textbar Ric\textbar\textless=K, its SAME bounded
near-closed form has L=3sqrt2, t=1/beta and avg
Z\_t\textless=C(r+D\_beta/V)+o(1). Sending the selectable allowance to
zero yields

\[
 \boxed{\frac{\operatorname{Vol}\{x:\operatorname{Vol}B_1(x)\ge v\}}V
       \le C(K,v)(r+D_\beta/V).}                      \tag{8.5}
\]

No class-square or Hessian budget is needed for (8.5). Hence any
bounded-Ricci sequence with actual compactified-pencil defect tending to
zero must collapse in measure at unit scale, for every fixed
v\textgreater0. A single noncollapsed point is not excluded when total
volume tends to infinity; the assertion concerns its volume fraction.
Uniform unit-ball noncollapse, or a fixed positive fraction of
noncollapsed points, gives a uniform defect gap.

\hypertarget{a-diameter-consequence-for-the-original-pointwise-ricci-input}{%
\subsubsection{9. A diameter consequence for the original pointwise
Ricci
input}\label{a-diameter-consequence-for-the-original-pointwise-ricci-input}}

Let
delta=\textbar\textbar Ric\^{}sharp+I\textbar\textbar{}\emph{\{infinity,op\},
diam(g)\textless=D, and suppose the actual class has p\textgreater0 and
a\^{}2\textgreater=T}+. There is epsilon(D)\textgreater0 such that

\[
                  \boxed{\delta^2+D_\beta/V\ge\epsilon(D),
                              \qquad\beta\ge4.}       \tag{9.1}
\]

To prove this it suffices to consider delta\textless=1/4. Then
\textbar Ric\textbar\textless=5/2, r\textless=4delta\^{}2, and
Chern--Weil gives

\[
 T_+\ge\frac{V}{4\pi^2}
       \left[\frac23(1-\delta)^2-2\delta^2\right]
       \ge V/(16\pi^2)>0.
\]

Integrality gives T\_+\textgreater=1, hence a\^{}2\textgreater=1 and
p\textgreater=1. The actual scalar bound and
\textbar s\textbar\textless=4(1+delta) give
V\textgreater=2pi\^{}2/(1+delta)\^{}2. Ricci lower bound and bounded
diameter then imply Vol B\_1(x)\textgreater=v(D)\textgreater0 at every
x. Apply (8.5), and absorb the remaining delta\textgreater=1/4 case by
decreasing epsilon(D) below 1/16.

Alternatively, replace the square condition by an explicit
p\textgreater=A\_0\textgreater0; the same proof gives epsilon(D,A\_0).
It is the noncollapse produced by these inputs that matters. The result
excludes actual near-saturation in this range; it neither constructs
Einstein metrics nor proves the unrestricted integral endpoint gap. It
does not assume asphericity, vanishing pi\_2, bilateral saturation, or
an opposite-Weyl energy bound.

\hypertarget{a-nonvanishing-upper-energy-margin-at-the-zero-weyl-corner}{%
\subsubsection{10. A nonvanishing upper-energy margin at the zero-Weyl
corner}\label{a-nonvanishing-upper-energy-margin-at-the-zero-weyl-corner}}

Consider an exact metric Ric=-g satisfying the actual hypotheses of
Section 9, with diam(g)\textless=D. Put T\_+=A(1-h),
p\textgreater=A\textgreater0 and e=R/(8pi\^{}2 A). Then V=2pi\^{}2 A e.
Write z=p/A\textgreater=1 and d=3(1-h)-e\textgreater=0; the last
inequality is the positive-chirality Chern--Weil identity. Direct
substitution gives

\[
 \frac{D_\beta}{V}=16+\frac{16}{\beta}-\frac{36z}{e}
                         +\frac{2(\beta+1)d}{e}.       \tag{10.1}
\]

Let epsilon=epsilon(D) in (9.1), decreased to at most one. Set

\[
 \beta_0=32/\epsilon,\quad a_0=\epsilon-16/\beta_0=\epsilon/2,
 \quad H=2(\beta_0-7)+a_0,\quad
 \gamma=\frac{9a_0}{4H}>0.
\]

\textbf{Corollary 8.} All these exact candidates obey

\[
 \boxed{e\le U_D(h):=
       \frac{6(\beta_0+1)(1-h)-36}{H}
       =\frac94-\gamma-
          \frac{6(\beta_0+1)}H(h-\tfrac14).}           \tag{10.2}
\]

Indeed D\_beta0/V\textgreater=epsilon by (9.1); use z\textgreater=1 in
(10.1), multiply by e, substitute d, and rearrange. H is positive. In
particular h\textgreater=1/4 forces e\textless=9/4-gamma. Let

\[
 G(h)=\frac{(4+\sqrt{1-3h})^2}{9},\qquad
 \eta_D=\min\{1/24,\gamma/10\}.
\]

Since G(1/4)=9/4 and \textbar G'\textbar\textless=5 on {[}1/4,7/24{]},
(10.2) gives

\[
 \boxed{e\le G(h)-\gamma/2,
              \qquad 1/4\le h\le1/4+\eta_D.}          \tag{10.3}
\]

The positive margin does not vanish as h approaches 1/4. It depends on
D; it is not an all-diameter improvement of the earlier upper envelope.
The alternative p\textgreater=A\_0 hypothesis gives the corresponding
D,A\_0 dependence.

There is also a diameter-free conclusion about any hypothetical sequence
at this corner. Suppose exact negative Einstein metrics with actual
p\_j\textgreater=A\_j satisfy h\_j-\textgreater1/4 and
e\_j-\textgreater9/4. Then d\_j-\textgreater0. Choose
beta\_j-\textgreater infinity so slowly that beta\_j d\_j-\textgreater0.
Equation (10.1), z\_j\textgreater=1, and D\_beta\textgreater=0 imply
D\_beta\_j/V\_j-\textgreater0. Applying (8.5) proves, for every fixed
v\textgreater0,

\[
 \frac{\operatorname{Vol}\{x:\operatorname{Vol}B_1(x)\ge v\}}{V_j}
                        \longrightarrow0.            \tag{10.4}
\]

Thus approaching this energy corner would force collapse in measure, not
merely unbounded diameter. This is a necessary degeneration statement,
not the construction of a sequence realizing the corner.

\hypertarget{local-exponential-averaging-excluding-bounded-curvature-collapse}{%
\subsubsection{11. Local exponential averaging: excluding
bounded-curvature
collapse}\label{local-exponential-averaging-excluding-bounded-curvature-collapse}}

The collapse conclusion in Section 8 need not be the endpoint. A second
averaging argument removes the noncollapse condition when sectional
curvature is bounded. Crucially, it does not assume that an exponential
chart is a regular covering or has uniformly bounded multiplicity.

Fix K,L\textgreater0 and a radius r\_0\textgreater0 with
4r\_0\textless pi/(2sqrt(max(K,1))). On a closed oriented four-manifold
let omega and t satisfy the field conditions of (8.1). Define the
curvature-regular set

\[
 G_{K,r_0}=\{x: |\sec_g|\le K\text{ on }B_{4r_0}(x)\}.
\]

\textbf{Theorem 9.} Without any curvature assumption outside these
balls,

\[
 \boxed{\frac{\operatorname{Vol}(G_{K,r_0})}{V}
                   \le C(K,L,r_0)\fint Z_t.}           \tag{11.1}
\]

In particular, if \textbar sec\textbar\textless=K globally then avg
Z\_t\textgreater=c(K,L)\textgreater0. For the actual conformal monopole
supplier with p\textgreater0, this yields

\[
 \boxed{r_{\rm Ric}+D_\beta/V\ge c(K)>0,
                 \qquad |\sec|\le K,\quad\beta\ge4.}  \tag{11.2}
\]

There is no diameter, injectivity-radius, volume or class-square
condition in (11.2). The extra curvature hypothesis is explicit; (11.2)
does not replace the unrestricted integral-Ricci problem.

\textbf{Proof.} For x in G, pull back the metric and form by exp\_x to
the tangent ball of radius 4r\_0. The map is nonsingular by the
conjugate-radius estimate. The pullback metric has controlled sectional
curvature and interior injectivity, hence a fixed-size harmonic chart on
a smaller ball. These normal-cover facts require no curvature
derivatives or injectivity on M; one precise formulation is
\href{https://arxiv.org/html/2601.16176v1\#S2.SS3}{Deleporte--Lagace--Rouveyrol,
Proposition 2.6}. Its proof is local, so bounds on B\_\{4r\_0\}(x)
suffice.

The contradiction proof of Theorem 6 on such a chart, with fixed r\_0,
gives epsilon\_loc(K,L,r\_0)\textgreater0 such that

\[
 E_{r_0}(x):=\int_{|v|<r_0\subset T_xM}
                Z_t(\exp_xv)\,d\mu_{\exp_x^*g}(v)
                              \ge\epsilon_{\rm loc}. \tag{11.3}
\]

This uses the exact same lifted self-dual form. If (11.3) failed, the
interior harmonic chart would give the forbidden smooth system (8.3).
The local cover is used only for that compactness argument, not to
impose any new gauge solvability assumption.

Let J\_x(s,theta) be its polar volume density. Rauch comparison gives
J\_x(s,theta)\textless=C(K,r\_0)s\^{}3 for 0\textless s\textless r\_0.
Therefore

\[
\begin{aligned}
 \epsilon_{\rm loc}\operatorname{Vol}(G)
 &\le C(K,r_0)\int_0^{r_0}s^3
          \int_G\int_{S_xM}Z_t(\exp_x(s\theta))d\theta d\mu(x)ds\\
 &\le C(K,r_0)|S^3|\frac{r_0^4}{4}\int_M Z_t.
\end{aligned}                                                     \tag{11.4}
\]

For the second line first drop the nonnegative center restriction G. The
geodesic flow preserves Liouville measure on the unit tangent bundle, so
for each s the integral over all (x,theta) is
\textbar S\^{}3\textbar int\_M Z\_t. This identity is valid even when
geodesics close or exp\_x has many sheets. It avoids both a false
uniform multiplicity bound and any deck-normality requirement. Equations
(11.1)--(11.2) follow, sending the actual supplier's selectable
allowance to zero. QED.

Define the sectional-curvature radius

\[
 r_{\sec}(x)=\sup\{0<r\le1:\sup_{B_r(x)}|\sec|\le r^{-2}\}.
\]

For any fixed a\textgreater0 the set \{r\_sec\textgreater=a\} is
contained in one G\_\{K,r\_0\}, with K,r\_0 depending only on a; choose
a smaller radius if the supremum is not attained. Consequently an actual
sequence with r\_Ric+D\_beta/V-\textgreater0 and beta\textgreater=4
satisfies

\[
       \frac{\operatorname{Vol}\{r_{\sec}\ge a\}}V\to0
                                  \quad\text{for every }a>0.    \tag{11.5}
\]

This forces curvature scale to vanish in measure, not merely injectivity
scale. It does not say pointwise curvature tends to infinity everywhere:
a high-curvature set can have large metric neighborhoods. In particular,
the exact energy-corner sequence in (10.4), if it existed, would also
obey (11.5). Section 10's nonvanishing energy-wall margin holds with a
K-dependent constant under \textbar sec\textbar\textless=K, without its
diameter or square hypothesis, by substituting (11.2) for (9.1).

\hypertarget{exact-einstein-metrics-replacing-curvature-bounds-by-euler-density}{%
\subsubsection{12. Exact Einstein metrics: replacing curvature bounds by
Euler
density}\label{exact-einstein-metrics-replacing-curvature-bounds-by-euler-density}}

The preceding argument can be continued once more. For exact Einstein
metrics, the curvature-regular fraction follows from an integral
curvature budget, even when balls collapse. This uses the unnormalized
small-energy theorem of Cheeger--Tian, not a regularity assertion for
arbitrary small Ricci residuals. The theorem is stated explicitly in
\href{https://sites.duke.edu/scshgap/files/2016/08/dim4dukebeamer.pdf}{Cheeger's
author presentation, pages 2--3} of the
\href{https://doi.org/10.1090/S0894-0347-05-00511-4}{2006 paper}.

Use the Hilbert--Schmidt norm of the full curvature operator R\_op on
Lambda\^{}2; fixed norm factors can be absorbed into the classical
constants. For Ric=-g, Chern--Gauss--Bonnet gives

\[
          \int|R_{\rm op}|^2=8\pi^2\chi(M),\qquad
          \Lambda=8\pi^2\chi(M)/V.                    \tag{12.1}
\]

\textbf{Theorem 10.} Given Lambda\_0,L\textgreater0, every closed Ric=-g
four-manifold with Lambda\textless=Lambda\_0, and every bounded
self-dual form and t in {[}0,1/4{]} as in (8.1), satisfies

\[
                         \fint Z_t\ge c(\Lambda_0,L)>0.
                                                               \tag{12.2}
\]

For actual conformal monopole input with p\textgreater0 this gives

\[
       \boxed{D_\beta/V\ge c(\Lambda_0)>0,
                       \qquad \beta\ge4.}            \tag{12.3}
\]

No diameter, total-volume, noncollapse, derivative, or class-square
bound is imposed. Here the Einstein equation is exact.

\textbf{Proof.} Cheeger--Tian supply constants epsilon\_CT,C\_CT such
that for 0\textless R\textless=1,

\[
 \int_{B_R(x)}|R_{\rm op}|^2\le\epsilon_{\rm CT}
        \quad\Longrightarrow\quad
 \sup_{B_{R/2}(x)}|R_{\rm op}|\le C_{\rm CT}R^{-2}.    \tag{12.4}
\]

This implication has no lower ball-volume premise. Put
F\_R(x)=int\_\{B\_R(x)\}\textbar R\_op\textbar\^{}2. Bishop comparison
for Ric=-g and Fubini give, with a dimensional B and R\textless=1,

\[
 \int_M F_R(x)d\mu(x)
 =\int_M|R_{\rm op}(y)|^2\operatorname{Vol}B_R(y)d\mu(y)
                         \le B R^4\Lambda V.          \tag{12.5}
\]

Choose R=min\{1,{[}epsilon\_CT/(2B Lambda\_0){]}\^{}\{1/4\}\}. At least
half the volume then consists of centers with F\_R\textless=epsilon\_CT.
On each corresponding B\_\{R/2\} the sectional curvature is bounded by
C\_CT R\^{}\{-2\}. Choose a fixed small theta\_0\textgreater0, depending
only on C\_CT, and put r\_0=theta\_0 R, K=max(C\_CT,1)R\^{}\{-2\}.
Choose theta\_0 so that 4r\_0\textless R/2 and
4r\_0\textless pi/(2sqrt(K)). All these centers belong to G\_\{K,r\_0\}.
Theorem 9 therefore yields 1/2\textless=C(K,L,r\_0)avg Z\_t, proving
(12.2). Apply the same actual form and send its selectable allowance to
zero to obtain (12.3). QED.

The proof also gives, for each fixed 0\textless R\textless=1, a useful
tradeoff

\[
           1\le C_0\Lambda R^4+C_1(R,L)\fint Z_t.      \tag{12.6}
\]

Indeed the high-energy centers have fraction at most C\_0 Lambda R\^{}4,
and all remaining centers are in the curvature-regular set controlled by
Theorem 9. Thus any exact Einstein actual zero-end sequence must have
chi(M\_j)/V\_j-\textgreater infinity. Curvature scale tending to zero is
not a free escape at bounded Euler density.

For the topological energy variables, e=V/(2pi\^{}2
A)\textgreater=p/A\textgreater=1 by the scalar monopole inequality.
Therefore

\[
                   \Lambda=4\frac{\chi/A}{e}.
\]

An independently checkable budget chi/A\textless=B\_0 supplies
Lambda\textless=4B\_0. The energy-wall proof of Section 10 consequently
holds with an epsilon depending only on B\_0, and with no diameter
assumption. In particular the corner h-\textgreater1/4,
e-\textgreater9/4 cannot be approached by these exact candidates.

This budget is automatic in the complex-hyperbolic surgery family: if X
is a closed complex-hyperbolic surface, A=c\_1\^{}2(X)=3chi(X), and
M=X\#k CPbar\^{}2\#ell(S\^{}1 x S\^{}3), then

\[
 \chi(M)/A=1/3+k/A-2\ell/A\le1/3+h,
                    \qquad h=(k+4\ell)/A.             \tag{12.7}
\]

Thus for 1/4\textless=h\textless=1/3 the density threshold is uniform
over all such actual inputs, their covers and all normalized exact
Einstein candidates. This excludes an energy region for arbitrary exact
outputs on the smooth types, not just for a chosen gluing family. It
does not claim that the entire remaining energy interval is realizable
or that these topologies are all non-Einstein solely from the upper
wall.

\hypertarget{a-uniform-forbidden-band-for-approximate-pointwise-inputs}{%
\subsubsection{13. A uniform forbidden band for approximate pointwise
inputs}\label{a-uniform-forbidden-band-for-approximate-pointwise-inputs}}

The energy-wall conclusion can be made robust without discarding scalar
variance. In this section retain the bounded-diameter hypotheses of
Section 9, with p\textgreater=A\textgreater0, T\_+=A(1-h), and
1/4\textless=h\textless=1/3. Either a\^{}2\textgreater=T\_+ or the
alternative p\textgreater=A\_0 is required. Let epsilon\_* denote the
resulting constant in (9.1), decreased to at most 1/16.

\textbf{Theorem 11.} There is c(D)\textgreater0 (or
c(D,A\_0)\textgreater0 for the alternative input) such that

\[
 \boxed{\delta_-^2+(G(h)-e)_+\ge c(D),
                       \qquad 1/4\le h\le1/3.}        \tag{13.1}
\]

One may take c=min\{1/64,epsilon\_*\^{}2/80000\}. In particular this is
a uniform forbidden neighborhood of the zero-Weyl upper-energy corner
for approximately Einstein input metrics, not just for exact outputs.

\textbf{Proof.} Put J=4P/R, ell=R/(4V), z=p/A-1\textgreater=0, B=3(1-h),
and b=B/e-1+J=12 int\textbar W\^{}+\textbar\^{}2/R\textgreater=0. Exact
Chern--Weil gives

\[
 \frac{D_\beta}{V}
  =\ell\left[16-4J-\frac{36(1+z)}e+2b+2\beta b
                                      +\frac{16-4J}{\beta}\right].
                                                               \tag{13.2}
\]

The scalar mean and variance are retained through

\[
 \ell=\frac{\bar s^2+L_s/V}{16(1-J/4)}
      =1-\tfrac12\fint(s+4)+\tfrac14r_{\rm Ric}.
\]

Let s\_0=(G(h)-e)\emph{+. If delta}-\textgreater1/8 or
s\_0\textgreater1/16, (13.1) follows immediately with the stated smaller
c.~Otherwise

\[
 J\le4\delta_-^2/(1-\delta_-)^2\le4/49,\quad
 247/144\le e\le49/20,\quad
 3/2\le e_0:=e(1-J)\le9/4.
\]

Here the upper bound on e follows from positive-chirality Chern--Weil,
e(1-J)\textless=B, and G\textgreater=16/9 gives its lower bound. Set
b\_0=B/e\_0-1. Then 0\textless=b\_0\textless=1/2 and
b=(1-J)b\_0\textless=(4-J)/8. Minimizing (13.2) over beta\textgreater=4
is therefore permitted at beta=sqrt(2(4-J)/b); for b=0 take the limit
beta-\textgreater infinity. The result is ell F\_\{J,z\}, where

\[
 F_{J,z}=16-4J-36(1+z)/e+2b+4\sqrt{2(4-J)b}.
\]

Define, for 3/2\textless=v\textless=B,

\[
 F_0(h,v)=16-36/v+2(B/v-1)+8\sqrt{2(B/v-1)}.
\]

Direct expansion yields

\[
 F_{J,0}-F_0(h,e_0)
 =J(-4+36/e_0-2b_0)
       +4\sqrt{2b_0}\,[\sqrt{(4-J)(1-J)}-2]\le20J.    \tag{13.3}
\]

Also vF\_0={[}4sqrt(v)+sqrt(2(B-v)){]}\^{}2-36. Its upper zero is G(h),
F\_0\textless=0 for v\textgreater=G(h), and elementary square-root
comparison gives F\_0(h,v)\textless=30sqrt((G(h)-v)\emph{+) on the
displayed compact v-range. Since (G-e\_0)}+\textless=s\_0+(49/20)J,
ell\textless=81/64, ell/e\textgreater=5/16, and
J\textless=256delta\_-\^{}2/49, (13.3) gives the convenient uniform
bound

\[
 \delta_-^2+\inf_{\beta\ge4}D_\beta/V
        \le200(\delta_-+\sqrt{s_0})-45z/4.
\]

The lower bound (9.1) holds for every beta, so
epsilon\_*+45z/4\textless=200(delta\_-+sqrt(s\_0)). Squaring and using
(delta\_-+sqrt(s\_0))\^{}2\textless=2(delta\_-\^{}2+s\_0) proves the
claim. QED.

For exact Einstein metrics this proof also applies with epsilon\_* from
Theorem 10 instead of (9.1). Thus chi/A\textless=B\_0, actual
p\textgreater=A, and 1/4\textless=h\textless=1/3 imply

\[
                      e\le G(h)-c(B_0),                \tag{13.4}
\]

without a diameter or class-square bound. In the complex-hyperbolic
surgery family (12.7), B\_0=2/3 suffices. The approximate-input version
of (13.4) does not follow from Cheeger--Tian: exact Einstein regularity
was essential in Theorem 10. These are distinct quantified extensions,
not a single unconditional estimate.

For exact metrics a shorter optimization gives the better explicit
margin c(B\_0)=epsilon\^{}2/288, where epsilon=min\{1,c(4B\_0)\} from
(12.3). Put a=3(1-h), d=a-e, and Delta=G(h)-e. If
beta\_*=sqrt(8e/d)\textgreater=4, minimizing (10.1), using
p/A\textgreater=1, gives

\[
 \epsilon e\le 14e+2a-36+8\sqrt{2e(a-e)}=:P_a(e).
\]

The upper zero of P\_a is G. Thus e\textless=G, and
e(a-e)\textless=G(a-G)+G Delta implies P\_a(e)\textless=8sqrt(2G Delta)
\textless=12sqrt(2Delta). Since e\textgreater=1,
Delta\textgreater=epsilon\^{}2/288. If beta\_*\textless4 then
e\textless2a/3 and G-2a/3\textgreater=4/9, which is stronger. The d=0
case follows by taking beta-\textgreater infinity. Therefore

\[
                 e\le G(h)-\epsilon^2/288.             \tag{13.5}
\]

The main improvement over the inherited upper wall is uniformity at
h=1/4, not the first appearance of a strict upper wall for fixed
h\textgreater1/4. As a secondary recovery of the previously studied
near-1/3 obstruction, choose
beta\_c=(14-epsilon+sqrt((14-epsilon)\^{}2+128))/4 in {[}7,8{]}. Its
e-independent term in (10.1) equals epsilon, so the remaining numerator
must be nonnegative. This gives
h\textless=1-6/(beta\_c+1)\textless=1/3-8epsilon/243. It is not counted
as a separate new obstruction mechanism.

\hypertarget{what-the-upper-wall-does-and-does-not-decide-topologically}{%
\paragraph{13.1 What the upper wall does and does not decide
topologically}\label{what-the-upper-wall-does-and-does-not-decide-topologically}}

The new fixed margin excludes candidate energy levels, not automatically
new smooth topologies. The recorded Chern--Weil budgets give upper
bounds e\textless=U, not a lower pinning of e near G. This matters for
the realized families. For instance take complex-hyperbolic covers with
A\_j-\textgreater infinity, let k\_j=ceil(A\_j/4), and make no S1xS3
surgery. The actual exceptional-sign classes give p\textgreater=A\_j and
a\^{}2=T\_+=A\_j-k\_j, while

\[
 h_j\ge1/4,\quad h_j\to1/4,\qquad
 \chi(M_j)/A_j\to7/12,\qquad U_j=3(1-h_j)\to9/4.
\]

At the limiting topological data, the formal candidate e=7/4, p/A=1 has
W\^{}+/(2pi\^{}2 A)=1/6 and W\^{}-/(2pi\^{}2 A)=1, where W\^{}pm here
denote their integrated squared energies. The scalar budget and every
classical pencil inequality are positive at this value; it stays
distance 1/2 below G(1/4). Thus these known inequalities do not force
candidates into the new forbidden upper strip. This is a witness about
implication between inequalities, NOT an Einstein metric or a
counterexample to a theorem.

More generally the known lower roots and the I41 wall remain separated
from G near h=1/4. After the harmless caps on the recorded constants, on
{[}1/4,7/24{]} their maximum is at most 19/12, whereas
G-1/288\textgreater7/4. Replacing the old upper root by (13.5) therefore
need not make the numeric feasible interval empty. An additional genuine
lower-volume/energy input would be needed for that use. In contrast, the
BF family in Section 4 already has an actual all-metric exclusion from
the LOWER linear wall; that application does not depend on such an
unproved upper-corner pinning.

\hypertarget{the-torsion-transition-cannot-be-closed-by-a-dirac-null-form-shortcut}{%
\subsubsection{14. The torsion transition cannot be closed by a Dirac
null-form
shortcut}\label{the-torsion-transition-cannot-be-closed-by-a-dirac-null-form-shortcut}}

This section records a targeted construction, not an actual gauge
counterexample. It explains why the signed transition flux in Theorem 5
still needs the full curvature equation or a global argument. The simple
pointwise Dirac constraints do not force its positive quartic component
to cancel.

Use Clifford matrices mapping S\^{}+ to S\^{}-:

\[
 C_1=I,\quad C_2=\operatorname{diag}(i,-i),\quad
 C_3=\begin{pmatrix}0&1\\-1&0\end{pmatrix},\quad
 C_4=\begin{pmatrix}0&-i\\-i&0\end{pmatrix}.
\]

Reverse Clifford action is -C\_j\^{}*. At a point Phi=sqrt(rho)e\_1,
write H\_\{ji\}=(h\_\{ji\},k\_\{ji\}) and K\_\{ij\}=k\_\{ij\}-k\_\{ji\}.
Differentiated D Phi=0 has components

\[
 d_{1j}=h_{j1}+ih_{j2}+k_{j3}-ik_{j4}=0,\quad
 d_{2j}=-h_{j3}-ih_{j4}+k_{j1}-ik_{j2}=0.              \tag{14.1}
\]

The linear Hessian contribution to nabla\_j t is
rho(k\_\{j1\}+ik\_\{j2\},-k\_\{j3\}-ik\_\{j4\}). The projected
derivative N\_\{bc\}=sum\_j(C\_j\^{}*)\_\{1b\}(nabla\_j t)\_c/sqrt(2)
can be written as

\[
 N=\sqrt{\rho/2}\,\mathcal L(\nabla\zeta)+\mathcal R_N,
 \quad |\mathcal R_N|\le C[\rho^{3/2}(|\operatorname{Ric}^{\circ}|
                        +|W^+-W_0|)+|\zeta||P_{\rm spin}|],      \tag{14.2}
\]

where

\[
\mathcal L(Z)=\begin{pmatrix}
 Z_{13}-Z_{24}+i(Z_{14}+Z_{23})&-Z_{33}+Z_{44}-i(Z_{34}+Z_{43})\\
 Z_{11}-Z_{22}+i(Z_{12}+Z_{21})&-Z_{31}+Z_{42}-i(Z_{32}+Z_{41})
\end{pmatrix}.
\]

For verification the Hessian expression is rho(L(h)+C(K))/sqrt(2),
modulo (14.1), with

\[
\mathcal C(K)=\begin{pmatrix}
2iK_{12}&-K_{13}-iK_{14}+iK_{23}-K_{24}\\
K_{13}+iK_{23}-iK_{14}+K_{24}&2iK_{34}
\end{pmatrix}.
\]

Central curvature, scalar curvature and the axial W\_0 model cancel in
this projection. Substitute sqrt(rho)h\_ji=nabla\_j
zeta\_i-\textless P\_j,P\_i\textgreater; the remaining quadratic
expression vanishes for the pure kernel component (iz,z,iw,w), while the
other terms are bounded by
C\textbar zeta\textbar\textbar P\_spin\textbar. This proves (14.2)
directly from the actual prolonged equations, not by discarding all
components except the Dirac trace.

In a compatible frame with t=sqrt(m)e\_1, writing x\_j=(nabla\_j t)\_1
gives the first column of N as (x\_1-ix\_2,-x\_3+ix\_4)/sqrt(2), while
dm=2sqrt(m) Re x. Hence

\[
                 |\langle t,\nabla t\rangle|^2
                      \le\tfrac34|dm|^2+4m|N|^2.      \tag{14.3}
\]

Determinant-curvature saturation does not directly control N, since that
curvature has canceled from (14.2). More decisively, even N=0 does not
remove the radial derivative in (14.3). Take h=0 and the symmetric
trace-free Hessian increment

\[
k=\frac12\begin{pmatrix}
1&-i&0&0\\-i&-1&0&0\\0&0&-1&i\\0&0&i&1
\end{pmatrix}.
\]

It satisfies (14.1), preserves curvature commutators and trace, and
gives nabla t/rho={[}(1,0),(-i,0),(0,1),(0,-i){]}. Its real Jacobian
determinant is positive one while N=0. These are exact Clifford-algebra
checks.

There is an actual local Dirac field realizing this increment. On flat
C\^{}2 with A=0, put z\_1=x\_1+ix\_2, z\_2=x\_3+ix\_4, g\_0=(bar
z\_1\^{}2-bar z\_2\^{}2)/4, and

\[
 \Phi_\epsilon(x)=(1,\epsilon g_0(x/\epsilon)),\qquad
 t_\epsilon(x)=(\bar z_1/\epsilon,\bar z_2/\epsilon).
                                                               \tag{14.4}
\]

Direct differentiation gives D\^{}+Phi\_epsilon=0 and
D\^{}-t\_epsilon=0, and t\_epsilon is exactly the cubic torsion of
Phi\_epsilon, not an independently chosen field. On
\textbar x\textbar\textless=R\_0 epsilon, for fixed R\_0,
rho=1+O(epsilon\^{}2), \textbar zeta\textbar=O(epsilon), \textbar d
zeta\textbar=O(1), and int\textbar nabla t\textbar\^{}2=O(epsilon\^{}2).
Yet its isolated torsion zero has positive index and can carry unit
signed transition flux.

Indeed at t=(r,0), write nabla t=(x+iy,z+iw) as real one-forms. In the
flat rank-two connection the transgression in Section 7 gives

\[
        d\chi(m)\wedge T_t
             =-\frac{\chi'(m)}{\pi^2m}
                       x\wedge y\wedge z\wedge w.      \tag{14.5}
\]

For (14.4) the real Jacobian is epsilon\^{}\{-4\} and
m=\textbar x\textbar\^{}2/epsilon\^{}2. If chi equals one near zero and
vanishes before m=R\_0\^{}2, polar integration gives int dchi wedge
T\_t=-int\_0\^{}\{R\_0\^{}2\}chi'(m)dm=1. Thus small current, the two
exact Dirac equations, and a small integrated first-derivative budget
cannot alone cancel the signed flux.

This field is NOT an actual positive-weight monopole: F\_A\^{}+=0 while
f sigma(Phi\_epsilon) is nonzero for f\textgreater0. Nor does flat
curvature satisfy the normalized negative Einstein target on a global
fixed-volume manifold. It rules out precisely the proposed
Dirac-only/null-form deduction, not the zero-end gap and not the full
actual equations. It leaves the new geometric exclusions in Sections
8--13 intact.

\hypertarget{the-positive-transition-core-persists-for-the-actual-local-gauge-equation}{%
\subsubsection{15. The positive transition core persists for the actual
local gauge
equation}\label{the-positive-transition-core-persists-for-the-actual-local-gauge-equation}}

The preceding construction can be continued to solve the positive-weight
Seiberg--Witten equations themselves. This strengthens the local test:
the missing global estimate cannot be obtained merely by declaring that
the critical transition scaling violates the local gauge PDE.

\textbf{Theorem 12 (local actual core).} Fix f\_*\textgreater0 and
a\textgreater0. For every sufficiently small epsilon\textgreater0 there
is a smooth unitary determinant connection A\_epsilon and positive
spinor Phi\_epsilon on the Euclidean ball B\_\{2epsilon\} such that

\[
 D_{A_\epsilon}\Phi_\epsilon=0,\qquad
       -iF_{A_\epsilon}^+=f_*\sigma(\Phi_\epsilon).
                                                               \tag{15.1}
\]

On smaller fixed-proportion balls their cubic torsion t\_epsilon has
exactly one positive-index zero, at the center. There is a cutoff
chi(\textbar t\_epsilon\textbar\^{}2), one near its zero and zero near
the boundary, for which

\[
 \int d\chi\wedge T_{t_\epsilon}=1+O(\epsilon^4).      \tag{15.2}
\]

The physical-coordinate estimates are

\[
\begin{gathered}
 |\Phi_\epsilon|^2=a^2+O(\epsilon^2),\quad
 |\langle\Phi_\epsilon,\nabla_A\Phi_\epsilon\rangle|=O(\epsilon),
 \quad |\nabla_A\Phi_\epsilon|=O(1),\quad |F_A|=O(1),\\
 \int|\nabla_A^2\Phi_\epsilon|^2+
       \int|\nabla_{3A}t_\epsilon|^2=O(\epsilon^2),\quad
 \int|F_A|^2=O(\epsilon^4).
\end{gathered}                                                     \tag{15.3}
\]

The signed transition flux nevertheless has order one, as does its
corresponding quartic transition energy. All these are unnormalized
local integrals. The flat core also has
int\textbar Ric+g\textbar\^{}2=O(epsilon\^{}4).

\textbf{Proof of existence.} Work first on the fixed Euclidean ball B\_2
in y-coordinates. Let lambda=epsilon\^{}2 and Phi\_0=a e\_1+epsilon
g\_0(y)e\_2, with g\_0 from (14.4). Thus D\^{}+Phi\_0=0. On
imaginary-valued forms put T=(d\^{}\emph{,sqrt(2)d\^{}+). In flat space
TT\^{}}=Delta on Omega\^{}0 direct-sum Omega\^{}2\_+, with Delta
nonnegative. Let G\_D be the componentwise Dirichlet inverse of Delta on
this smooth ball, and define

\[
 R_A=T^*G_D,\qquad R_D=D^-G_D.
\]

Then TR\_A=Id and D\^{}+R\_D=Id. On C\^{}\{k,alpha\} for fixed
k\textgreater=3, 0\textless alpha\textless1, both are bounded operators
gaining one derivative, by the classical
\href{https://ocw.mit.edu/courses/18-156-differential-analysis-spring-2004/pages/lecture-notes/}{Dirichlet
Schauder estimates, lectures 11--17}. Their boundary traces are not
prescribed independently; only their right-inverse property is used. For
a spinor Phi in a fixed bounded C\^{}\{k,alpha\} neighborhood of Phi\_0
set

\[
 A(\Phi)=R_A(0,\sqrt2\lambda i f_*\sigma(\Phi)),\qquad
 \mathcal T(\Phi)=\Phi_0-\tfrac12R_D[c(A(\Phi))\Phi]. \tag{15.4}
\]

Products and the quadratic spinor map are locally Lipschitz in this
Banach space. A(Phi)=O(lambda) and the Lipschitz constant of T is
O(lambda), uniformly for small epsilon. A ball of radius C lambda about
Phi\_0 is mapped to itself for fixed sufficiently large C; for
sufficiently small epsilon, T is a contraction. Its fixed point and the
associated A satisfy

\[
 d^*A=0,\quad F_A^+=\lambda i f_*\sigma(\Phi),\quad D_A\Phi=0,
 \quad \|A\|_{C^{k+1,\alpha}}+
        \|\Phi-\Phi_0\|_{C^{k,\alpha}}\le C\epsilon^2.
\]

The fixed-point identities and elliptic regularity bootstrap to
smoothness. Rescale x=epsilon y, setting Phi\_x(x)=Phi\_y(x/epsilon) and
A\_x(x)=epsilon\^{}\{-1\}A\_y(x/epsilon) in components. The Dirac
equation scales by epsilon\^{}\{-1\}, the curvature by
epsilon\^{}\{-2\}; lambda=epsilon\^{}2 cancels the latter factor. This
proves (15.1) with the fixed positive f\_*.

\textbf{Torsion and charge.} In y-coordinates the physical torsion obeys

\[
        t_\epsilon(\epsilon y)
           =a^2(\bar z_1,\bar z_2)+O(\epsilon)
                                      \quad\text{in }C^1(B_1).
                                                               \tag{15.5}
\]

This follows directly from the cubic formula: its leading derivative of
epsilon g\_0 cancels the physical epsilon\^{}\{-1\}, and all fixed-point
and connection corrections are O(epsilon). The leading real linear map
is invertible with positive determinant a\^{}8. Small C1 perturbation
makes the map injective on a smaller convex ball. The degree on its
boundary is positive one, so it has exactly one zero. In fact Phi\_0 is
even under y-\textgreater-y. The Dirichlet inverse preserves parity,
whereas T\^{}* and D\^{}- reverse it. Thus the contraction preserves
even Phi and produces odd A; the fixed point has these parities. Its
torsion is odd, placing its unique zero exactly at the center. The
bounds (15.3) follow by differentiating (15.4) and rescaling. In
particular the squared Hessian L2 integral is O(epsilon\^{}2); its L2
norm is O(epsilon), not O(epsilon\^{}2).

In the extending flat spin frame, E\_-=S\^{}- tensor L has curvature
(3/2)F\_A I and c\_2(E\_-)=-9F\_A wedge F\_A/(16pi\^{}2). Removing a
small sphere about the torsion zero, the transgression convention of
Section 7 gives

\[
 \int_B c_2(E_-)-\int_{\partial B}T_t=1,
 \qquad
 \int_B d\chi\wedge T_t=1-\int_B\chi c_2(E_-).
\]

Since F\_A=O(1) and Vol(B)=O(epsilon\^{}4), (15.2) follows. On a fixed
torsion-amplitude transition annulus, (15.5) also gives the order-one
quartic cost directly. QED.

The costs can be identified more precisely. The constant-source part of
the Dirichlet calculation gives, in y-components,

\[
 A_y=\epsilon^2 A_0+O(\epsilon^3),\qquad
 A_0=\frac{if_*}{2}\iota_y\sigma(ae_1).
\]

Consequently the physical curvature is if\_*sigma(ae\_1)+O(epsilon), its
ASD part is O(epsilon), and

\[
 \int_{B_{2\epsilon}}|F_A|^2
    =\pi^2f_*^2a^4\epsilon^4+O(\epsilon^5).
                                                               \tag{15.6}
\]

For fixed 0\textless m\_1\textless m\_2\textless4a\^{}4 let
T\_epsilon=\{m\_1\textless\textbar t\_epsilon\textbar\^{}2\textless m\_2\}.
The C1 convergence in (15.5), valid on B\_2, and rescaling imply

\[
\begin{aligned}
 \int_{T_\epsilon}|\nabla t|^4
       &\longrightarrow8\pi^2(m_2^2-m_1^2),\\
 \int_{T_\epsilon}|\langle t,\nabla t\rangle|^2|\nabla t|^2
       &\longrightarrow\frac{8\pi^2}{3}(m_2^3-m_1^3),\\
 \int_{T_\epsilon}|dm|^2|\nabla t|^2
       &\longrightarrow\frac{16\pi^2}{3}(m_2^3-m_1^3).
\end{aligned}                                                   \tag{15.7}
\]

Indeed the limiting y-field has m=a\^{}4\textbar y\textbar\^{}2,
\textbar dt\textbar\^{}2=4a\^{}4,
\textbar\textless t,dt\textgreater\textbar\^{}2=2a\^{}8\textbar y\textbar\^{}2
and \textbar dm\textbar\^{}2=4a\^{}8\textbar y\textbar\^{}2; integrate
with the four-dimensional polar factor 2pi\^{}2 r\^{}3 dr. The limiting
level spheres are regular, so replacing the exact transition region by
its limiting annulus is justified by this C1 convergence.

The leading anti-complex-linear derivative has N=0 in Section 14's
projection. The actual Phi-adapted frame differs by O(epsilon), while
the physical derivative is O(epsilon\^{}\{-1\}); hence N=O(1) and

\[
       \int_{T_\epsilon}m|N|^2|\nabla t|^2=O(\epsilon^2).
                                                               \tag{15.8}
\]

Thus the positive charge survives while this projected contribution
tends to zero. The critical radial mass term, not an uncontrolled N
term, can carry the transition cost. Even
int\textbar dm\textbar\^{}2=O(epsilon\^{}2) does not remove it.

For the flat metric the inherited majorant weight is exactly sqrt(3), so
f\_*=sqrt(3) is an allowed concrete choice; there is no zero-weight
limit. Choosing a\^{}2=4/3 and radius R\_0 epsilon with R\_0=2sqrt(2)/3
makes the leading boundary torsion mass a\^{}4 R\_0\^{}2=128/81 equal to
the zero-end mass. This is amplitude matching only. Its outgoing normal
mass derivative is 2a\^{}4 R\_0/epsilon+O(1), and its unit torsion
boundary map has degree +1 in the extending trivialization. It cannot be
joined through a zero-free collar to a regular nonzero torsion field
extending over a ball, whose boundary degree is zero. Compensating zeros
or a nontrivial marked bundle identification would be necessary.

No such exterior, clutching, fixed closed manifold, or conformally
solvable global class system is produced here. The flat scalar curvature
and weight also differ from the negative Einstein exterior model.
Consequently this is an exact local concentration mechanism with its
relative charge, not a counterexample to O10, not a closed near-Einstein
construction, and not a reversal of the geometric exclusions above. It
does establish that the remaining work must use global matching or
information beyond the local weighted gauge equations and the displayed
integral budgets alone.

\hypertarget{removing-the-local-metric-and-weight-mismatch}{%
\paragraph{15.1 Removing the local metric and weight
mismatch}\label{removing-the-local-metric-and-weight-mismatch}}

The local core also exists on a FIXED hyperbolic background with
sectional curvature -1/3. It then satisfies Ric=-g, W\^{}pm=0 and the
canonical weight f\_0=2 exactly; the two mismatches just noted are not
local obstructions. Here is the perturbation argument, including its
boundary convention.

Fix a point p.~In geodesic coordinates x=epsilon y set
g\_epsilon=epsilon\^{}\{-2\}(exp\_p(epsilon dot))\^{}*g on B\_2. In
polar and Cartesian coordinates, respectively,

\[
\begin{aligned}
 g_\epsilon&=dr^2+\frac3{\epsilon^2}
                    \sinh^2(\epsilon r/\sqrt3)g_{S^3},\\
 (g_\epsilon)_{ij}
  &=\delta_{ij}+\frac{\epsilon^2}{9}
             (|y|^2\delta_{ij}-y_i y_j)+O(\epsilon^4).
\end{aligned}                                                     \tag{15.9}
\]

In every fixed C\^{}\{k,alpha\} norm this is an O(epsilon\^{}2)
perturbation of the Euclidean metric. Use the positive-square-root
orthonormal coframe and its spin lift. The spin connections and
coefficients of D\_epsilon-D\_0 are O(epsilon\^{}2), the latter as a
first-order operator. For T\_epsilon=(d*\_epsilon,sqrt(2)d+\_epsilon)
the exact identity T\_epsilon T\_epsilon\^{}*=Delta\_H remains valid on
Omega\^{}0 direct-sum Lambda+\_epsilon. Identify that target with a
fixed trivial bundle. Both this second-order operator and D+\_epsilon
D-\_epsilon have the form

\[
          L_\epsilon=\Delta_0+O(\epsilon^2)\partial^2
                             +O(\epsilon^2)\partial+O(\epsilon^2).
\]

Their auxiliary Dirichlet inverses exist by the convergent Neumann
series

\[
 G_\epsilon=G_0[I+(L_\epsilon-\Delta_0)G_0]^{-1}.
\]

The norm here is C\^{}\{k,alpha\}-\textgreater C\^{}\{k,alpha\} for the
bracketed operator; G\_0 maps into C\^{}\{k+2,alpha\} with zero boundary
trace. Thus the new right inverses R\_A=T\_epsilon\^{}*G\_A and
R\_D=D-\_epsilon G\_D gain one derivative uniformly. No boundary
condition on A or Phi is asserted.

The flat seed Phi\_0=a e\_1+epsilon g\_0 e\_2 is no longer exactly
harmonic, but the corrected seed

\[
       \widehat\Phi_0=\Phi_0-R_D D_\epsilon^+\Phi_0
\]

is: D+\emph{epsilon Phi\_hat\_0=0, and
Phi\_hat\_0-Phi\_0=O(epsilon\^{}2). Apply the same contraction (15.4)
with this seed, the curved right inverses, lambda=epsilon\^{}2 and
f}*=2. The result obeys the exact physical equations

\[
            D_{A,g}\Phi=0,\qquad F_A^+=2i\sigma_g(\Phi).
                                                               \tag{15.10}
\]

The physical metric is the original fixed g, not g\_epsilon. It has
s=-4, W\^{}pm=0 and q=4, so the inherited majorant is indeed f\_0=2. All
fields differ from the flat seed in rescaled coordinates by
O(epsilon\^{}2). The physical cubic torsion therefore still has the C1
expansion (15.5), giving one positive simple zero and the same limiting
transition costs (15.7)--(15.8). Only its O(epsilon\^{}2) displacement
in physical coordinates is needed here; no parity assertion is
necessary.

Take a\^{}2=4/3. The physical determinant curvature has
F\_A\^{}-=O(epsilon) and
\textbar F\_A\^{}+\textbar\^{}2=8/9+O(epsilon\^{}2). The full bundle
E\_-=S\^{}- tensor L now includes the nonflat spin connection. Its
characteristic density is

\[
 16\pi^2c_2(E_-)
   =\left[9(|F_A^+|^2-|F_A^-|^2)+\frac23\right]d\mu_g
   =\left[\frac{26}{3}+O(\epsilon^2)\right]d\mu_g.
\]

It is bounded; the relative formula still gives int dchi wedge
T\_t=1+O(epsilon\^{}4). The Ricci and Weyl error terms vanish
identically on the core. The form-closure term does NOT vanish
identically, and no normalized full defect on these shrinking balls is
claimed to tend to zero. Quadratic derivative costs remain
O(epsilon\^{}2), whereas the critical radial transition cost has a
positive limit.

This removes the metric and canonical-weight mismatches, but not the
boundary winding, O(epsilon\^{}\{-1\}) outgoing torsion derivative,
determinant ASD matching, or global conformally solvable class. In
particular, the core's F\_A\^{}- is small while the globally calibrated
endpoint exterior requires its squared mean to approach 2/9. Exact local
Einstein and gauge equations admit the core; producing a closed endpoint
degeneration still requires the entire exterior and its matching.
Theorem 10 excludes that global construction at bounded Euler density,
so it is not contradicted.

END COMPLETE SOURCE

\end{document}
